\section{Notation and the basic package}\label{sec:basic}

Throughout, $G$ is a finite simple connected graph on $n \ge 2$ vertices with vertex set
$V$ and edge set $E$, $|E| = m$. We abbreviate $\alpha = \alpha(G)$, $f = f(G)$,
$d = \diam(G)$, $\Delta = \Delta(G)$, $\nabla = \nabla(G) = n - f$, $\mu = \mu(G)$. For
$S \subseteq V$, $G[S]$ is the induced subgraph, $N(v)$ the neighbourhood of $v$ and
$N[v] = N(v) \cup \{v\}$.

\subsection{The Havel--Hakimi trajectory}

Run the Havel--Hakimi reduction on the degree sequence of $G$ with vertex labels attached,
so that at each step we may speak of \emph{which vertex} is deleted. Let $s = s(G)$ be the
number of steps performed before the process halts, let $D_1, \dots, D_s$ be the values of
the successive deleted entries (the \emph{heads}), let $K$ be the set of vertices deleted
as heads, and call the $n - s$ vertices never deleted the \emph{survivors}. The
\emph{block} of step $j$, written $\mathrm{block}_j$, is the set of $D_j$ entries
decremented at that step: the prefix of the remaining entries, sorted non-increasingly,
after deleting the step-$j$ head.

Ties are not broken canonically by the definition, and the campaign's protocol treats
tie-sensitivity as a first-class failure mode: every numerical run cited below is performed
under the canonical sort \emph{and} under several randomised adversarial tie-breaks, and
every proof step below that could depend on a tie is either tie-invariant or says which
tie-break it fixes and why fixing it is legitimate.

\begin{lemma}[$\res = n - s$]\label{lem:hh}
\Selem
Let $s(G)$ be the number of Havel--Hakimi steps and $D_1,\dots,D_s$ the successive head
values. Then
\begin{enumerate}[label=\textup{(\arabic*)}]
\item $\res(G) = n - s(G)$;
\item $D_1 = \Delta \ge D_2 \ge \cdots \ge D_s \ge 1$ and $\sum_{i=1}^{s} D_i = m$.
\end{enumerate}
\end{lemma}

\begin{proof}
(1) Each step shortens the list by exactly one entry, and the recursion stops on a list
whose head --- hence, the list being sorted non-increasingly, every entry --- is $0$,
returning its length. After $s$ steps the length is $n - s$.

(2) Havel--Hakimi preserves graphicality \cite{havel,hakimi}, so at every step the head
$D_i$ is at most the number of remaining strictly positive entries; the $\N$-truncated
subtraction therefore never truncates, and one step decreases the sum of the list by
exactly $2D_i$ ($D_i$ removed as the head, $D_i$ further units removed by the decrements).
The sum starts at $2m$ and ends at $0$, giving $\sum_i D_i = m$. For monotonicity, the new
list is obtained from the old one by deleting the head and decreasing some entries, so its
maximum cannot increase.
\end{proof}

\begin{fact}[Favaron--Mah\'eo--Sacl\'e]\label{fact:fms}
\Scomp{asserted by the verification script on all $30\,995$ graphs of stage E and never
violated; cited, not reproved}
$\res(G) \le \alpha(G)$ for every graph $G$.
\end{fact}

This is used as a black box throughout. It is due to Favaron, Mah\'eo and Sacl\'e
\cite{fms}; a short proof is due to Griggs and Kleitman \cite{gk}. The bibliographic data
was supplied by an automated literature pass and should be re-checked before any
formalization (Section~\ref{sec:limits}); the inequality itself is asserted by the
campaign's verification script on the whole exhaustive range and never fails there.

\subsection{The cover \texorpdfstring{$B$}{B} and the reductio hypothesis}
\label{sec:cover}

Fix a \emph{maximum} independent set $A$ and put $B := V \setminus A$, $\tau := |B| = n -
\alpha$. Then $B$ is a minimum vertex cover, every $a \in A$ has $N(a) \subseteq B$ and
hence $\dg(a) \le \tau$, and every $b \in B$ has at least one neighbour in $A$ (otherwise
$A \cup \{b\}$ would be independent). We write $\degA(v) = |N(v) \cap A|$,
$\degB(v) = |N(v) \cap B|$, $e_B = |E(G[B])|$, and
\[
\nu := \binom{\tau}{2} - e_B
\]
for the number of non-adjacent pairs inside $B$. Since $A$ is independent, every edge has
an endpoint in $B$, so $m = \sum_{b \in B} \dg(b) - e_B$.

By Lemma~\ref{lem:hh}(1) and Fact~\ref{fact:fms}, $n - s = \res \le \alpha = n - \tau$, so
$s \ge \tau$ always. The campaign's central standing hypothesis is the opposite extreme:

\begin{definition}[reductio hypothesis]\label{def:reductio}
The \emph{reductio hypothesis} is $\res(G) = \alpha(G)$, equivalently $s(G) = \tau(G)$.
\end{definition}

Under the reductio hypothesis the Havel--Hakimi process must terminate in exactly $\tau$
steps, and it is that scarcity of steps which the whole of
Sections~\ref{sec:toolkit}--\ref{sec:gfannu} exploits. We also fix the low/high split of
$B$ that organises everything from Section~\ref{sec:budget} onwards:
\[
\Bhi := \{b \in B : \dg(b) \ge \tau+1\}, \quad p := |\Bhi|, \qquad
\Blo := B \setminus \Bhi, \quad L := |\Blo| = \tau - p.
\]
For $S \subseteq B$ let $\nu(S)$ be the number of non-adjacent pairs inside $S$, and put
$\mbar := \nu(\Bhi)$. A vertex $b \in B$ is \emph{$B$-universal} if it is adjacent to every
other vertex of $B$; $\Bloplus$ denotes the set of low vertices that are \emph{not}
$B$-universal.

\section{Elementary certificates, and the concentration at diameter four}
\label{sec:elementary}

This section is self-contained and elementary. Every proof in it is short and is given in
full, so that the reader can discharge the whole section without reference to the review
process; only Fact~\ref{fact:fms} and the classical graphicality theorem behind
Lemma~\ref{lem:hh} are used as black boxes.

\subsection{Diameter at most three}

\begin{lemma}\label{lem:f-alpha}
\Selem
For every connected $G$ with $n \ge 2$: $f(G) \ge \alpha(G) + 1$.
\end{lemma}

\begin{proof}
$G$ is connected with $n \ge 2$, so it has an edge; hence no independent set is all of $V$.
Take a maximum independent set $A$ and any $v \in V \setminus A$. In $G[A \cup \{v\}]$ the
set $A$ is independent, so every edge is incident with $v$: the graph is the star with
centre $v$ and leaves $N(v) \cap A$, together with $|A| - |N(v) \cap A|$ isolated vertices.
A star plus isolated vertices is acyclic, and $|A \cup \{v\}| = \alpha(G) + 1$.
\end{proof}

\begin{theorem}\label{thm:diam3}
\Selem
If $\diam(G) \le 3$ then \eqref{eq:C61} holds.
\end{theorem}

\begin{proof}
$n \ge 2$ and $G$ connected give $d \ge 1$, so $\lceil d/3 \rceil = 1$. By
Fact~\ref{fact:fms} and Lemma~\ref{lem:f-alpha},
$\res(G) + 1 \le \alpha(G) + 1 \le f(G)$.
\end{proof}

This alone settles $89.7\%$ of connected graphs on at most seven vertices and $78.4\%$ of
those on eight.

\subsection{The 2-packing certificate}

Call $B' \subseteq V$ a \emph{2-packing} if $\dist(b,b') \ge 3$ for all distinct
$b,b' \in B'$.

\begin{lemma}[star-forest lemma]\label{lem:starforest}
\Selem
If $A$ is independent, $B'$ is a 2-packing and $A \cap B' = \emptyset$, then
$G[A \cup B']$ is a forest --- indeed a disjoint union of stars and isolated vertices.
\end{lemma}

\begin{proof}
$B'$ is independent because its members are at distance $\ge 3 > 1$, and $A$ is
independent; hence every edge of $G[A \cup B']$ joins $A$ to $B'$. If distinct
$b,b' \in B'$ had a common neighbour $a \in A$ then $\dist(b,b') \le 2$, a contradiction;
so the sets $N(b) \cap A$, $b \in B'$, are pairwise disjoint. Therefore $G[A \cup B']$ is
the disjoint union of the stars with centre $b$ and leaf set $N(b) \cap A$, together with
the uncovered vertices of $A$.
\end{proof}

\begin{theorem}[2-packing certificate]\label{thm:packing}
\Selem
Let
$\rho^{*}(G) := \max\{ |B'| : B'$ a 2-packing with $B' \cap A = \emptyset$ for some maximum
independent set $A \}$.
Then $f(G) \ge \alpha(G) + \rho^{*}(G)$; consequently, if
$\rho^{*}(G) \ge \lceil d/3 \rceil$ then \eqref{eq:C61} holds for $G$.
\end{theorem}

\begin{proof}
Lemma~\ref{lem:starforest} applied to a witnessing pair $(A,B')$ gives an induced forest of
order $\alpha + |B'|$; then apply Fact~\ref{fact:fms}.
\end{proof}

\begin{corollary}\label{cor:B1}
\Selem
For every connected $G$ with $n \ge 2$:
$f(G) \ge \alpha(G) + \lfloor (d-1)/4 \rfloor + 1$.
\end{corollary}

\begin{proof}
Fix a maximum independent set $A$ and a shortest path $P = v_0v_1\cdots v_d$ realising the
diameter; a shortest path is induced and $\dist(v_i,v_j) = |i-j|$. Since $A \cap V(P)$ is
independent in $P$, no two consecutive indices lie in $A$; hence
$I := \{ i \in [0,d] : v_i \notin A \}$ meets $\{i,i+1\}$ for every $0 \le i \le d-1$.
Consequently consecutive elements of $I$ differ by at most $2$, $\min I \le 1$ and
$\max I \ge d-1$.

Greedily set $i_1 := \min I$ and $i_{j+1} := \min\{ i \in I : i \ge i_j + 3 \}$ while such
an $i$ exists. Because consecutive elements of $I$ differ by at most $2$ we have
$i_{j+1} \le i_j + 4$, so $i_k \le 1 + 4(k-1)$. The process stops at index $k$ with
$\max I < i_k + 3$, i.e.\ $i_k > \max I - 3 \ge d-4$. Combining,
$1 + 4(k-1) > d-4$, i.e.\ $k > (d-1)/4$, so $k \ge \lfloor (d-1)/4 \rfloor + 1$.

The set $B' := \{v_{i_1},\dots,v_{i_k}\}$ has pairwise index gaps $\ge 3$, hence pairwise
distances $\ge 3$ (equalities hold along a shortest path), and $B' \cap A = \emptyset$ by
construction. Apply Lemma~\ref{lem:starforest}.
\end{proof}

\begin{corollary}\label{cor:B2}
\Selem
\eqref{eq:C61} holds for every connected $G$ with
$\diam(G) \in \{1,2,3,5,6,9\}$.
\end{corollary}

\begin{proof}
By Corollary~\ref{cor:B1} and Fact~\ref{fact:fms} it suffices that
$\lfloor (d-1)/4 \rfloor + 1 \ge \lceil d/3 \rceil$. For $d \ge 10$ this fails, since
$\lfloor (d-1)/4 \rfloor + 1 \le (d+3)/4 < d/3 \le \lceil d/3 \rceil$, the middle
inequality being $3d + 9 < 4d$. A finite check of $d = 1,\dots,9$ leaves exactly
$d \in \{1,2,3,5,6,9\}$.
\end{proof}

The constant $1/4$ here, rather than the $1/3$ the conjecture asks for, is exactly the
quantitative gap; see Obstacle~O2 in Section~\ref{sec:obstacles}.

\subsection{The matching certificate}

\begin{theorem}\label{thm:matching}
\Selem
If $\mu(G) \ge \nabla(G) + \lceil d/3 \rceil$ then \eqref{eq:C61} holds.
\end{theorem}

\begin{proof}
A maximum matching has $\mu$ pairwise disjoint edges and an independent set contains at
most one endpoint of each, so $\alpha + \mu \le n$. Hence by Fact~\ref{fact:fms},
$\res(G) \le \alpha \le n - \mu \le n - \nabla - \lceil d/3 \rceil
= f(G) - \lceil d/3 \rceil$.
\end{proof}

\begin{corollary}\label{cor:C1}
\Selem
\eqref{eq:C61} holds whenever $\nabla(G) \le \lceil d/2 \rceil - \lceil d/3 \rceil$; in
particular for every forest, and for every graph of decycling number $0$.
\end{corollary}

\begin{proof}
A shortest path with $d$ edges contains a matching of size $\lceil d/2 \rceil$, so
$\mu(G) \ge \lceil d/2 \rceil$; apply Theorem~\ref{thm:matching}. For a forest,
$\nabla = 0 \le \lceil d/2 \rceil - \lceil d/3 \rceil$, the right-hand side being
non-negative for every $d$.
\end{proof}

\subsection{An equivalent form}

\begin{theorem}\label{thm:equivform}
\Selem
\eqref{eq:C61} is equivalent to
\[
s(G) \;\ge\; \nabla(G) + \left\lceil \frac{\diam(G)}{3} \right\rceil,
\]
that is: \emph{the Havel--Hakimi process must run for at least $\lceil d/3 \rceil$ steps
more than the decycling number}.
\end{theorem}

\begin{proof}
Immediate from $\res = n - s$ (Lemma~\ref{lem:hh}) and $f = n - \nabla$.
\end{proof}

The classically known part of this is $s \ge \nabla + 1$, which is Fact~\ref{fact:fms}
combined with Lemma~\ref{lem:f-alpha}. So the conjecture asks for a gain of
$\lceil d/3 \rceil - 1$ over the classical bound, and \emph{only} that.

\subsection{Falsification, and where the residual difficulty sits}
\label{sec:falsify}

\begin{remark}[falsification stage]\label{rem:falsify}
\Scomp{$150\,905$ graph evaluations in stages A--D, $0$ violations; stage E re-checks
$30\,995$ of them against three further invariants, also $0$ violations;
\texttt{wowii61\_verify.py all}, runtime $\approx 100$\,s on one core}
No counterexample to \eqref{eq:C61} was found. The stages are: (A) every connected graph on
$2 \le n \le 7$ vertices, $995$ graphs, from the graph atlas; (B) a complete cover of all
$11\,117$ connected graphs on $8$ vertices, obtained by enumerating all
$853 \times 127 = 108\,331$ pairs (connected $H$ on $7$ vertices, non-empty
$N(v) \subseteq V(H)$) --- every connected graph has a non-cut vertex, so every connected
$8$-vertex graph is such an $H + v$, and no isomorphism rejection is needed since duplicates
only cost time; (C) $1\,209$ structured graphs with $n \le 26$ (paths, cycles, path and
cycle powers $P_n^k, C_n^k$ for $k \le 4$, lexicographic blow-ups $P_m[K_r]$, $C_m[K_r]$,
$P_m[E_r]$, $C_m[E_r]$, chains of cliques sharing cut vertices, theta graphs, caterpillars,
random trees); (D1) $7\,994$ uniform random connected graphs on $n = 9,10$; (D2) $32\,376$
graphs on $n = 9,\dots,13$ produced by simulated annealing that minimises the slack
$f - \res - \lceil d/3 \rceil$. Stage E additionally \emph{asserts} $\res \le \alpha$,
$f \ge \alpha+1$, $f \ge \alpha + \lfloor (d-1)/4 \rfloor + 1$ and \eqref{eq:C61} itself on
every graph it touches, so a counterexample to any of them aborts the run.
\end{remark}

The minimum of the slack $f - \res - \lceil d/3 \rceil$ over all connected graphs with
$n \le 8$, by diameter, is
\[
\begin{array}{c|ccccccc}
d & 1 & 2 & 3 & 4 & 5 & 6 & 7\\ \hline
\min \text{slack} & 0 & 0 & 0 & \mathbf{0} & 1 & 2 & 2\\
\min (f-\alpha) & 1 & 1 & 1 & \mathbf{1} & 2 & 2 & 4
\end{array}
\]
so in that range the inequality is tight exactly up to $d = 4$ and becomes loose beyond.

\begin{remark}[measured coverage]\label{rem:coverage}
\Scomp{all $995$ connected graphs with $n \le 7$ and a $30\,000$-graph prefix of the $n=8$
enumeration}
Evaluating the three proved certificates P1 $=$ Theorem~\ref{thm:diam3},
P2 $=$ Theorem~\ref{thm:matching}, P3 $=$ Theorem~\ref{thm:packing} gives coverage
$89.7 / 74.1 / 96.1\%$ and $98.8\%$ jointly for $n \le 7$, and
$78.4 / 98.0 / 93.3\%$ and $99.2\%$ jointly at $n = 8$. Of the $267$ graphs in that range
covered by none of the three, \emph{every one has diameter exactly $4$}; among them
$f - \alpha = 1$ for $223$ and $= 2$ for $44$, $\alpha - \res \in \{1,2\}$, the slack is $0$
for $112$ and $1$ for $155$, and $\nabla \in \{1,2,3\}$.
\end{remark}

We stress the scope of Remark~\ref{rem:coverage}: it is a statement about $n \le 8$. The
diameters left open in general by Corollary~\ref{cor:B2} are $4$, $7$, $8$ and every
$d \ge 10$; what the measurement says is that in the exhaustively checked range the
residual instances all sit at $d = 4$, not that $d = 4$ is the only residual diameter for
all $n$. Section~\ref{sec:limits} returns to this.

\subsection{The diameter-four dichotomy}\label{sec:dichotomy}

At $d = 4$ we have $\lceil d/3 \rceil = 2$, while Corollary~\ref{cor:B1} yields only
$f \ge \alpha + 1$ and Fact~\ref{fact:fms} yields only $\res + 2 \le \alpha + 2$. The
shortfall is exactly one unit, which is the following statement.

\begin{proposition}[the $d=4$ dichotomy]\label{prop:dichotomy}
\Selem
Suppose that for every connected $G$ with $\diam(G) = 4$,
\begin{equation}\label{eq:D4}
f(G) \ge \alpha(G) + 2 \quad\text{or}\quad \res(G) \le \alpha(G) - 1.
\tag{D4}
\end{equation}
Then \eqref{eq:C61} holds for every connected $G$ with $\diam(G) = 4$, and hence, with
Corollary~\ref{cor:B2}, for every connected $G$ with
$\diam(G) \in \{1,2,3,4,5,6,9\}$.
\end{proposition}

\begin{proof}
Let $\diam(G) = 4$, so $\lceil d/3 \rceil = 2$. In the first case
$\res + 2 \le \alpha + 2 \le f$ by Fact~\ref{fact:fms}. In the second,
$\res + 2 \le \alpha + 1 \le f$ by Lemma~\ref{lem:f-alpha}.
\end{proof}

Everything from Section~\ref{sec:toolkit} onwards is directed at \eqref{eq:D4}. Negating it
gives the object the campaign calls the hard core, in two tiers which must be kept apart ---
several statements below hold in the weaker one, and conflating them was itself a defect
found in review.

\begin{definition}[frame and hard core]\label{def:frame}
The \emph{hard-core frame} is: $G$ connected, not a forest, $\diam(G) = 4$, and
$f(G) = \alpha(G) + 1$. The \emph{hard core} is the frame together with the reductio
hypothesis $\res(G) = \alpha(G)$ --- and, after Theorem~\ref{thm:T3}, together with
$\tau \ge 4$.
\end{definition}

\eqref{eq:D4} is exactly the assertion that the hard core is empty. A third tier is needed
too, and is used for most of Section~\ref{sec:toolkit}: \emph{reductio only}, meaning the
hypothesis $\res = \alpha$ with no structural assumption whatever. That tier is not vacuous
even where the hard core is: $C_5$ has degree sequence $[2,2,2,2,2]$, $\alpha = 2$,
$\tau = 3$ and $\res = 2$, so it satisfies the reductio while $\diam = 2$ and
$f = 4 \ne \alpha+1$, i.e.\ it lies outside the frame entirely.

\subsection{Named obstacles}\label{sec:obstacles}

Four obstacles were isolated, each with an explicit witness. They explain why the
elementary package stops where it does, and they constrain what a proof can look like. All
witnesses below are re-verified by \texttt{wowii61\_verify.py w}.

\begin{description}
\item[O1: the two slacks do not separate.] Put $X := \alpha - \res \ge 0$ and
$Y := f - \alpha \ge 1$; then \eqref{eq:C61} says exactly $X + Y \ge \lceil d/3 \rceil$.
The tree $W_2$ on $8$ vertices with edges
$\{01,12,23,24,25,26,73\}$ has $d = 4$, $\alpha = 5$, $\res = 5$, $f = 8$, so $X = 0$ at
$d = 4$; the graph $W_1$ on $8$ vertices with edges
$\{01,12,13,14,25,56,71,75\}$ has $d = 4$, $\alpha = 6$, $f = 7$, $\res = 5$, so $Y = 1$ at
$d = 4$ (and $W_1$ is simultaneously a tight instance of \eqref{eq:C61}: $7 = 5+2$). Hence
no proof bounding one of $X$, $Y$ alone by $\lceil d/3 \rceil - O(1)$ can work: the
diameter information must enter the Havel--Hakimi process and the forest construction
simultaneously. In particular the two obvious strengthenings
``$\alpha + 1 \ge \res + \lceil d/3 \rceil$'' and ``$f \ge \alpha + \lceil d/3 \rceil$''
are both false, witnessed by $W_2$ and $W_1$ respectively.

\item[O2: the $1/4$ versus $1/3$ gap.] Corollary~\ref{cor:B1}'s bound
$\rho^{*} \ge \lfloor (d-1)/4 \rfloor + 1$ is best possible \emph{for that argument}: the
extremal configuration has $A \cap V(P) = \{v_0,v_2,v_4,\dots\}$, so the path vertices
available off $A$ are the odd-indexed ones and a 2-packing among them must skip every
second one. Getting $d/3$ requires admitting pairs $v_i, v_{i+2}$ into $B'$, which is
legitimate exactly when $v_i$ and $v_{i+2}$ have at most one common neighbour in $A$ ---
i.e.\ one must control $4$-cycles between consecutive odd-indexed path vertices and $A$.
The smallest graph where that control genuinely fails is $W_4$ on $6$ vertices with edges
$\{02,13,24,25,34,35\}$: $d = 4$, $\alpha = 4$, $f = 5$, $\res = 3$, slack $0$, and
$\rho^{*} = 1 < 2 = \lceil d/3 \rceil$.

\item[O3: $\res$ admits no usable vertex-deletion recursion.] Both natural recursions fail.
For $v$ of minimum degree, $\res(G) \le 1 + \res(G - N[v])$ is false: on $W_5$, $n=6$, edges
$\{02,05,13,14,23,34,35\}$, degrees $[2,2,2,4,2,2]$, $v = 0$, one has $\res(G) = 3$ but
$\res(G - N[0]) = 1$. For $v_0$ a diametral endpoint, $\res(G) \le \res(G - N[v_0])$ fails
--- for \emph{every} diametral endpoint --- in $36.4\%$ of the $7\,228$ connected graphs
with $n \le 8$ and $d \ge 4$. The structural reason is common to both: deleting $N[v]$
lowers $n$, which lowers $\res$, and lowers the degrees of the survivors, which raises it,
and the two effects are not comparable.

\item[O4: degree-sequence-only bounds are exactly one short at $d = 4$.] For $W_1$
($n=8$, $m=8$, $\Delta=5$, degrees $[5,3,2,2,1,1,1,1]$) every bound derivable from sequence
data alone gives $\res \le 6$: $\res \le \alpha = 6$; $\res \le n - \mu = 6$; and, from
Lemma~\ref{lem:hh}(2) with $D_i \le \Delta$, $\res \le n - \lceil m/\Delta \rceil = 6$. The
truth is $\res = 5$, and \eqref{eq:C61} needs exactly $\res \le f - 2 = 5$. So the residual
family demands a residue bound one better than every first-order degree-sequence estimate.
\end{description}

Two further routes were closed by proof rather than by witness. The ball-peeling
certificate $\Phi(G) := \max_{B'} (|B'| + \alpha(G - N[B']))$ over 2-packings $B'$ satisfies
$\Phi(G) \le \alpha(G)$ --- for a 2-packing $B'$ and an independent set $I$ of $G - N[B']$,
the set $B' \cup I$ is independent in $G$ --- so it can never beat Fact~\ref{fact:fms} and
dies exactly where O1 does; measured, it fails for $46.3\%$ of the connected graphs with
$n \le 8$, the identical count to $\alpha < \res + \lceil d/3 \rceil$. And replacing
$\alpha$ by $f$ in that recursion is vacuous, since
$\max_{B'}(|B'| + f(G - N[B'])) = f(G)$ attained at $B' = \emptyset$; so the layering idea
has to be grafted onto a maximum independent set, which is what
Theorem~\ref{thm:packing} does.

\section{The Havel--Hakimi toolkit}\label{sec:toolkit}

Everything in this section is proved under the reductio hypothesis alone, or with no
hypothesis at all. Lemmas~\ref{lem:S} and~\ref{lem:T} are two- and five-line arguments given
here in full and carry no registry row; the remaining statements are certified. Their
registry rows are R-1 to R-5 (and R-6,
R-7 for the two corollaries mentioned in Remark~\ref{sec:upstream-selfref}); the evidence chain is a
Qwen3.8-Max round, an independent Claude Opus 5 round, and a third, differently-briefed
Opus session that confirmed the repairs of the second, with the campaign's own
reproduction of every joint (Appendix~\ref{app:ledger}).

Recall the vocabulary of Section~\ref{sec:basic}: heads, survivors, blocks; and say an
entry is \emph{excess at step $j$} if its current value exceeds $s - j + 1$.

\begin{lemma}[survivor degree bound]\label{lem:S}
\Selem
Under the reductio hypothesis, every survivor has degree $\le \tau$; equivalently, every
vertex of degree $\ge \tau+1$ is a head, so $\#\{v : \dg(v) \ge \tau+1\} \le \tau$.
\end{lemma}

\begin{proof}
A survivor is decremented at most once per step and must reach $0$ after $\tau$ steps.
\end{proof}

\begin{lemma}[survivor decay, general $s$]\label{lem:F3prime}
\Scert{row R-5}
Let $s$ be the number of Havel--Hakimi steps. At the start of step $j$, $1 \le j \le s$,
every survivor has value $\le s - j + 1$. \textup{(}No reductio hypothesis.\textup{)}
\end{lemma}

\begin{proof}
After step $s$ every remaining entry is $0$ by Lemma~\ref{lem:hh}(1). A survivor is
decremented at most once per step and only the $s - j + 1$ steps $j, j+1, \dots, s$ remain.
\end{proof}

This lemma exists in the form stated because of a defect found in review. Its predecessor
was stated for $s = \tau$ only, yet was cited inside the proofs of
Lemmas~\ref{lem:Zplus} and~\ref{lem:DICH}, which are asserted for general $s$. The
statement of those two lemmas was true in the claimed generality all along; only the
citation reached outside the cited lemma's scope. Over $19\,324$ trajectories --- the
exhaustive connected atlas $n \le 7$ plus $4\,500$ random graphs $n = 8,\dots,12$, canonical
sort plus three adversarial tie-breaks --- of which $14\,784$ have $s \ne \tau$, i.e.\ lie
in exactly the regime where the old citation was invalid, Lemma~\ref{lem:F3prime} fails $0$
times and Lemma~\ref{lem:Zplus} fails $0$ times.

\begin{lemma}[terminal shape]\label{lem:T}
\Selem
$s = \tau$ if and only if the list at the start of step $\tau$ is exactly
$[D_\tau, 1^{D_\tau}, 0^{\dots}]$. Moreover, for an arbitrary graph
\textup{(}no reductio\textup{)}: if two entries are $\ge 2$ at the start of step $\tau$,
then $s \ne \tau$, whence by Fact~\ref{fact:fms} $s \ge \tau+1$ and so
$\res \le \alpha - 1$.
\end{lemma}

\begin{proof}
The last step deletes the head and decrements the next $D_\tau$ entries; all entries must
then be $0$, which is the displayed shape. For the second sentence, the displayed shape has
exactly one entry $\ge 2$ when $D_\tau \ge 2$ and none when $D_\tau = 1$; two entries $\ge
2$ therefore contradicts it, so $s \ne \tau$, and $s \ge \tau$ always by
Fact~\ref{fact:fms}.
\end{proof}

The second sentence is stated \emph{outside} the reductio deliberately. Read inside the
standing hypothesis $s = \tau$ its antecedent contradicts the first sentence and it is
vacuously true; the version above is the one every call site actually consumes. This
distinction was itself the subject of a review finding.

\begin{lemma}[block occupancy]\label{lem:Zplus}
\Scert{row R-1}
At the start of step $j$, every later head whose current value exceeds $s - j + 1$ lies in
$\mathrm{block}_j$, and is therefore decremented at step $j$.
\end{lemma}

\begin{proof}
Let $x$ be deleted at step $i > j$ with current value $v > s-j+1$, and suppose
$x \notin \mathrm{block}_j$. The block is a prefix of the remaining entries in
non-increasing order after deletion of the step-$j$ head, so $v \le \min(\mathrm{block}_j)$.
If some block entry is a survivor then $\min(\mathrm{block}_j) \le s-j+1$ by
Lemma~\ref{lem:F3prime}, contradicting $v > s-j+1$. Otherwise all $D_j$ block entries are
later heads; together with $x$ that is $D_j + 1$ later heads, while only $s - j$ heads
remain, so $D_j \le s-j-1$. But every block entry is at most the head value $D_j$, so
$v \le \min(\mathrm{block}_j) \le D_j \le s-j-1 < s-j+1$, a contradiction.
\end{proof}

\begin{lemma}[head dichotomy]\label{lem:DICH}
\Scert{row R-2}
Let a head be deleted at step $i$, of original degree $g$, and let $h_i$ be the number of
earlier steps at which it was decremented. Then:
\begin{enumerate}[label=\textup{(\alph*)}]
\item \emph{persistence}: if it is excess at step $j < i$ then by Lemma~\ref{lem:Zplus} it
is decremented at step $j$, hence it is excess at step $j+1$ as well;
\item if $g \ge s+1$ then it is excess at step $1$, hence at every step $1,\dots,i-1$, hence
in every earlier block, so $h_i = i-1$ and $D_i = g - (i-1)$ \emph{exactly};
\item if $g \le s$ then $D_i \le s - i + 2$, unconditionally.
\end{enumerate}
\end{lemma}

\begin{proof}
(a) and (b) are immediate from Lemma~\ref{lem:Zplus} and the definition of excess. For (c):
if the entry is never excess before deletion, its value at step $i-1$ is at most
$s - (i-1) + 1 = s-i+2$ and values never increase (for $i = 1$, $D_1 = g \le s$).
Otherwise let $j_0 \ge 2$ be the first excess step; its value at step $j_0-1$ is at most
$s - j_0 + 2$, and by (a) it is decremented at each of the $i - j_0$ steps
$j_0, \dots, i-1$, so $D_i \le (s-j_0+2) - (i-j_0) = s-i+2$.
\end{proof}

Lemma~\ref{lem:DICH} is the engine of everything downstream. Clause (b) is an
\emph{equality}, and it is what makes the position sums of Section~\ref{sec:budget}
collapse; clause (c) is an unconditional inequality for low heads. The pair replaces an
earlier, hypothesis-laden head-decay lemma of the campaign, and is tie-safe, which the
earlier one was not obviously.

\begin{theorem}[the $K = B$ corner is a clique]\label{thm:K}
\Scert{row R-3}
Assume the reductio hypothesis $\res(G) = \alpha(G)$ and that every $b \in B$ has
$\dg(b) \ge \tau+1$. Then $K = B$ and $e_B = \binom{\tau}{2}$: $B$ is a clique.
\end{theorem}

\begin{proof}
By Lemma~\ref{lem:S} and the hypothesis, every $b \in B$ is a head; since
$|K| = s = \tau = |B|$ we get $K = B$, and hence the total number of decrements absorbed by
heads before their deletion equals $e_B$. Every head now has original degree
$\ge \tau+1 = s+1$, so Lemma~\ref{lem:DICH}(b) applies to each, giving $h_i = i-1$.
Summing, $e_B = \sum_{i=1}^{\tau}(i-1) = \binom{\tau}{2}$.
\end{proof}

\begin{corollary}\label{cor:K1}
\Scert{row R-4}
Let $G$ be connected with $\diam(G) = 4$. If some maximum independent set $A$ has
$\min_{b \in B} \dg(b) \ge \tau+1$, then $\res(G) \le \alpha(G) - 1$.
\end{corollary}

\begin{proof}
Otherwise $\res = \alpha$, and Theorem~\ref{thm:K} makes $B$ a clique, contradicting
Observation~\ref{obs:R1} below.
\end{proof}

\begin{remark}[independent numerical check]\label{rem:thmK}
\Scomp{exhaustive $n \le 7$, the exhaustive $n = 8$ cover, and $6\,000$ denser random
graphs $n = 9,\dots,12$: $114\,914$ connected graphs, of which $1\,464$ have
$\res = \alpha$ with all $B$-degrees $\ge \tau+1$}
In none of those $1\,464$ instances is $B$ a non-clique. The check was written by the owner
agent and reuses no code from the agent that proposed Theorem~\ref{thm:K}.
\end{remark}

\begin{remark}[a statement of the campaign refuted by its own review]
\label{sec:upstream-selfref}
An earlier statement of this line asserted, unconditionally, that if every $b \in B$ has
$\dg(b) \ge \tau+1$ and $e_B = 0$ then $m \le \Delta(G) + \tau(\tau+1)/2 - 1$. A reviewing
session produced $K_{2,3}$ as a counterexample: taking $B$ to be the $2$-element part,
$\alpha = 3$, $\tau = 2$, $e_B = 0$, both $B$-degrees are $3 = \tau+1$, $\Delta = 3$ and
$m = 6$, while $\Delta + \tau(\tau+1)/2 - 1 = 5 < 6$. The owner agent reproduced the
witness and upheld the objection: the statement is true only under the reductio hypothesis
--- and indeed $K_{2,3}$ has heads $(3,2,1)$, so $s = 3 \ne \tau = 2$ and the reductio fails
there, which is precisely why the hypothesis has to be written into the statement. The
patched statement carries registry row R-7. Nothing downstream used the unconditional
version. We report this because the honest measure of a review gate is what it catches, and
because it is the smaller of the two refutations in this paper; the larger one is in
Section~\ref{sec:gfannu}.
\end{remark}

\section{Frame lemmas at diameter four}\label{sec:frame}

The four statements of this section hold in the hard-core frame
(Definition~\ref{def:frame}) --- no reductio hypothesis --- and are the only structural
input the later sections take from the geometry of $G$. Each has a short proof, given here
in full; none carries a registry row of its own.

\begin{lemma}[pairing obstacle]\label{lem:pair}
\Selem
Let $A$ be a maximum independent set, $B = V \setminus A$, and suppose $f = \alpha + 1$.
Then for distinct $u,v \in B$:
\[
uv \in E \implies N_A(u) \cap N_A(v) \ne \emptyset, \qquad
uv \notin E \implies |N_A(u) \cap N_A(v)| \ge 2 .
\]
\end{lemma}

\begin{proof}
Consider $G[A \cup \{u,v\}]$, which has $\alpha + 2$ vertices. If $f = \alpha+1$ it must
contain a cycle. Since $A$ is independent, every edge of that subgraph meets $\{u,v\}$, so a
cycle is either the triangle $u - v - a$ with $a \in N_A(u) \cap N_A(v)$, or the $4$-cycle
$u - a - v - b$ with distinct $a,b \in N_A(u) \cap N_A(v)$. The first requires $uv \in E$.
\end{proof}

\begin{observation}\label{obs:R1}
\Selem
If $\diam(G) = 4$ then $B$ is not a clique; equivalently $\nu \ge 1$.
\end{observation}

\begin{proof}
$A$ is a maximal independent set, so every $a \in A$ has a neighbour in $B$. If $B$ were a
clique then for $a, a' \in A$ we would have $a - b - b' - a'$ for suitable $b,b' \in B$, so
$\dist(a,a') \le 3$; for $a \in A$, $b \in B$ we would have $\dist(a,b) \le 2$; and within
$B$, $\dist \le 1$. Hence $\diam \le 3$.
\end{proof}

\begin{fact}[the diameter witness]\label{fact:Fb}
\Selem
Call $T \subseteq B$ an \emph{occurring type} if $T = N(a)$ for some $a \in A$. In the
hard-core frame there exist occurring types $T_1, T_2$ that are disjoint and have no
$G[B]$-edge between them.
\end{fact}

\begin{proof}
By Lemma~\ref{lem:pair}, any two vertices of $B$ have a common $A$-neighbour or are
adjacent, so $\dist(b,b') \le 2$ and hence $\dist(a,b) \le 3$ for all $a \in A$,
$b \in B$. So $\diam = 4$ can only be realised by a pair $a, a' \in A$, and
$\dist(a,a') = 4$ holds if and only if $N(a) \cap N(a') = \emptyset$ and there is no
$G[B]$-edge between them.
\end{proof}

\begin{lemma}[$\tau$-uniform degree bound on the witness types]\label{lem:Cstar}
\Selem
In the hard-core frame, let $T_1, T_2$ realise Fact~\ref{fact:Fb}. Then every
$b \in T_1 \cup T_2$ has $\degA(b) \ge 3$, hence $\dg(b) \ge 3 + \degB(b)$.
\end{lemma}

\begin{proof}
Let $b \in T_1$ and pick any $b' \in T_2$. There is no $G[B]$-edge between $T_1$ and $T_2$,
so $bb' \notin E$, and Lemma~\ref{lem:pair} gives $|N_A(b) \cap N_A(b')| \ge 2$. The witness
$a_1$ of type $T_1$ is adjacent to $b$ and, since $b' \notin T_1$, not adjacent to $b'$; so
$a_1$ is none of those two common neighbours. Hence $\degA(b) \ge 2 + 1 = 3$. Symmetrically
for $b \in T_2$.
\end{proof}

\begin{remark}[the measured obstacle]\label{rem:obstacle-tau}
\Scomp{exhaustive $n \le 7$ plus the $n = 8$ cover, restricted to the hard-core frame:
$332$ graphs, $399$ pairs realising Fact~\ref{fact:Fb}, $0$ failures of
Lemma~\ref{lem:Cstar}}
The frame instances split by $\tau$ as $\{2 : 14,\ 3 : 200,\ 4 : 118\}$, and the hypothesis
of Corollary~\ref{cor:K1} --- all $B$-degrees $\ge \tau+1$ --- holds in $14$ of $14$, $91$
of $200$ and only $5$ of $118$ of them respectively. Lemma~\ref{lem:Cstar} explains the
collapse: Fact~\ref{fact:Fb} together with Lemma~\ref{lem:pair} forces $\degA(b) \ge 3$
only for the $O(1)$ vertices of $T_1 \cup T_2$, and $\degA(b) \ge 2$ for any $b$ with a
non-neighbour in $B$ --- bounds that do not grow with $\tau$, while the hypothesis of
Corollary~\ref{cor:K1} demands $\tau+1$. This is why the $K = B$ corner is a shrinking
corner and why the rest of the paper is about $K \ne B$.
\end{remark}

\section{Small covers: \texorpdfstring{$\tau \le 3$}{tau <= 3}}\label{sec:smalltau}

\begin{theorem}[$\tau \le 2$]\label{thm:tau2}
\Selem
If $\diam(G) = 4$ and $\tau(G) = n - \alpha(G) \le 2$, then \eqref{eq:D4} holds for $G$.
\end{theorem}

\begin{proof}
We may assume $f = \alpha + 1$, since otherwise $f \ge \alpha+2$ and there is nothing to
prove; so we are in the hard-core frame and Lemma~\ref{lem:pair} is available.

If $\tau = 1$ then $B$ is a single vertex covering all edges, so $G$ is a star and
$\diam \le 2$. Let $\tau = 2$, $B = \{x,y\}$. Every $a \in A$ has $N(a) \subseteq \{x,y\}$
non-empty, so $A$ partitions into $A_x$ (type $\{x\}$), $A_y$ (type $\{y\}$) and $A_c$
(type $\{x,y\}$); write $a = |A_x|$, $b = |A_y|$, $c = |A_c|$. By Fact~\ref{fact:Fb} there
are two disjoint occurring types with no $G[B]$-edge between them, which forces
$xy \notin E$ and $A_x, A_y \ne \emptyset$; and then Lemma~\ref{lem:pair} applied to the
non-adjacent pair $x,y$ gives $|A_c| = c \ge 2$. Thus
\[
\dg(x) = a + c, \qquad \dg(y) = b + c, \qquad m = a + b + 2c ,
\]
and the degrees of the $A$-vertices are $1$ (on $A_x \cup A_y$) or $2$ (on $A_c$).

Suppose $\res = \alpha$, i.e.\ $s = \tau = 2$, so $D_1 + D_2 = m$ by
Lemma~\ref{lem:hh}(2). Assume without loss of generality $a \ge b$, so
$D_1 = \Delta = a+c$. After the first step the entry $y$ has value $\dg(y) - 1$ or
$\dg(y)$ according to whether $y \in \mathrm{block}_1$; but $\dg(y) = b + c \ge c \ge 2$ is
the largest of the remaining entries, so $y$ heads the sorted order and lies in
$\mathrm{block}_1$ as soon as $D_1 \ge 1$, which holds. Hence
$D_2 \le (b+c) - 1$ and
\[
D_1 + D_2 \;\le\; (a+c) + (b+c) - 1 \;=\; m - 1 \;<\; m,
\]
a contradiction. Therefore $s \ge 3$ and $\res = n - s \le n - 3 = \alpha - 1$.
\end{proof}

\begin{theorem}[$\tau = 3$]\label{thm:T3}
\Scert{row R-13, with repairs R1, R3, J1, J2, K1, K2, K3 landed}
If $G$ is connected, $\diam(G) = 4$ and $\tau(G) = 3$, then \eqref{eq:D4} holds for $G$.
Combined with Theorem~\ref{thm:tau2}, \eqref{eq:D4} holds whenever $\tau \le 3$; so the
hard core has $\tau \ge 4$.
\end{theorem}

The proof is a finite case analysis on
$k := \#\{b \in B : \dg(b) \ge 4\}$; since every $A$-vertex has degree $\le \tau = 3$, only
$B$-vertices can have degree $\ge 4$. It is long, and we give its skeleton together with
the tools it uses, rather than reproducing every branch; the complete branch-by-branch text,
with the seven landed repairs listed in the marker, is in the campaign's source report, and
it is that text --- not this summary --- that the two review rounds certified. We flag this
explicitly in Section~\ref{sec:limits} as the one place in the paper where a certified
argument is summarised rather than reproduced.

\begin{proof}[Proof skeleton]
Assume $f = \alpha+1$ and $\res = \alpha$, so $s = 3$, $\sum_{i \le 3} D_i = m$, and by
Lemma~\ref{lem:T} the list after step $2$ is exactly $[D_3, 1^{D_3}, 0^{\dots}]$. Write
$B = \{x,y,z\}$ with $B$-degrees $X \ge Y \ge Z$, and for non-empty $T \subseteq B$ let
$\mathrm{mu}[T] = \#\{a \in A : N(a) = T\}$; put $p = \mathrm{mu}[B]$,
$q = \sum_{|T| = 2}\mathrm{mu}[T]$, $r = \sum_{|T|=1}\mathrm{mu}[T]$, so
$\alpha = p+q+r$, $n = \alpha+3$, and the $A$-degree spectrum is $3^p 2^q 1^r$. Since $A$
is independent, $G$ is determined up to isomorphism by these seven multiplicities together
with the isomorphism type of $G[B]$, of which there are four; and $e_B \le 2$ by
Observation~\ref{obs:R1}. Three structural bounds follow from Lemma~\ref{lem:pair} and
Fact~\ref{fact:Fb}:
\begin{itemize}
\item if $e_B = 0$ then every $b \in B$ has $\dg(b) \ge 2$;
\item if $e_B = 1$, say $E(G[B]) = \{uv\}$ with third vertex $w$, then
$\dg(u),\dg(v),\dg(w) \ge 3$ and $k \ge 1$;
\item if $e_B = 2$, so $G[B]$ is a path with centre $c$ and ends $e_1,e_2$, then
$\dg(e_1),\dg(e_2) \ge 4$ --- hence $k \ge 2$ --- and $\dg(c) \ge 3$.
\end{itemize}
Consequently $k = 0 \Rightarrow e_B = 0$ and $k = 1 \Rightarrow e_B \le 1$. One further
tool is used: if the degree list with one copy of $\Delta$ removed has all entries $\le 3$
\emph{and} $\Delta$ is at least the number of its entries equal to $3$, then after step $1$
all entries are $\le 2$, so $D_2, D_3 \le 2$ and $\sum_{i \le 3} D_i \le \Delta + 4$; hence
$m > \Delta + 4$ forces $s \ge 4$. Both hypotheses of this tool must be checked at each
call site --- the second is not automatic, and the campaign's own control instance
$[3,3,3,3,3,2,1]$ (with $\Delta = 3$ and $m = 9 > 7$) is exactly where a blanket appeal to
it would give a false bound. That over-general appeal was a defect found in review and
removed; the tool is now invoked only at sites where its second hypothesis is verified in
line.

The four branches then run as follows. For $k = 3$ one gets $D_1 = X$, $D_2 = Y - 1$ and
$D_3 = Z + 1 - e_B$, and comparing with the maximum possible value after step $2$ forces
$e_B = 2$ and $D_3 = Z-1$ with $Z = 4$; counting the entries of value $\ge 2$ after step $2$
then gives at least $Y \ge 4$ of them, contradicting Lemma~\ref{lem:T}. For $k = 2$ the
three values of $e_B$ are treated separately: $e_B = 0$ forces $Z \le 2$ and an undecremented
$3$ at step $1$, whence $p \ge X \ge 4$ and so $Z \ge p \ge 4$, a contradiction;
$e_B = 1$ forces $\degA(z) = p = 3$ and $\mathrm{mu}[\{x,z\}] = \mathrm{mu}[\{y,z\}] =
\mathrm{mu}[\{z\}] = 0$, which kills every pair admissible for Fact~\ref{fact:Fb};
$e_B = 2$ pins the degree sequence to
$[a_0+c_0+2,\, b_0+c_0+2,\, 3,\, 3,\, 2^{c_0},\, 1^{a_0+b_0}]$ with $a_0, b_0, c_0 \ge 1$,
and the explicit trajectory gives $\sum_{i\le3} D_i = m - 1 < m$. For $k = 1$ the two
possible values of $e_B$ each reduce to $p \le 2 \le X - 2$, where the tool above applies
and contradicts the exact value of $m$; in the sub-branch $e_B = 0$, $p = 3$ every occurring
type contains one fixed vertex $x$, so every $A$-vertex is adjacent to $x$, every two
$A$-vertices are at distance $2$ through $x$, $\dist(a,b) \le 3$ and $\dist(b,b') = 2$,
whence $\diam \le 3$ unconditionally --- contradicting $\diam = 4$. Finally $k = 0$ forces
$e_B = 0$, then $p = 2$ and $q = 1$, and Fact~\ref{fact:Fb} pins the graph up to isomorphism
to the single instance $A = \{u_1,u_2 \text{ universal},\, w \sim \{y,z\},\, \ell \sim
\{x\}\}$, $\dg = [3,3,3,3,3,2,1]$, $n = 7$, $m = 9$; running the process directly gives
$D = (3,3,2)$, $\sum D_i = 8 = m-1 < m$, so $s = 4$ and $\res = 3 = \alpha - 1$. These
four cases are exhaustive, so $\res = \alpha$ is impossible.
\end{proof}

\begin{remark}[$\tau = 3$: numerical backing, and non-vacuity]\label{rem:T3num}
\Scomp{$591\,710$ hard-core-frame instances at $\tau = 3$ inside a multiplicity box of size
$5$, and $147\,348$ instances inside a box of size $4$ with the structural skeleton checked
separately}
The parameterisation above makes the $\tau = 3$ frame an exhaustive scan rather than a
search: every tuple in the box is covered, no isomorphism rejection is needed, and no SAT
solver is involved. In the size-$5$ box there are $0$ instances with $\res \ge \alpha$; the
distribution by $e_B$ is $\{0 : 224\,035,\ 1 : 196\,075,\ 2 : 171\,600\}$ with $e_B = 3$
never occurring, which is a numerical confirmation of Observation~\ref{obs:R1}; and the
minimum of $m - (D_1+D_2+D_3)$ is exactly $1$, attained at $B$ independent,
$A = \{u_1, u_2$ universal, $w \sim \{x,y\}$, $e \sim \{z\}\}$, $n = 7$, $m = 9$,
$\dg = [3,3,3,3,3,2,1]$. In the size-$4$ box the $(e_B,k)$ cells $(1,0)$, $(2,0)$ and
$(2,1)$ are always empty, confirming the two implications $e_B = 1 \Rightarrow k \ge 1$ and
$e_B = 2 \Rightarrow k \ge 2$ used above.
Independently, the round-B judge enumerated every $\tau = 3$ graph on $n \le 37$
($89.8$M parameter tuples), a further $77.9$M-tuple box up to $n = 73$, $7.7$M
huge-multiplicity instances with $n$ up to about $10^6$, and a $2^{21}$ brute force at
$n = 7$ reproducing exactly the five hard-core degree sequences its own parameterisation
predicts, with no counterexample anywhere --- and, decisively for the question of vacuity,
$24\,013$ near-misses: configurations satisfying every hard-core condition except
$f = \alpha+1$. Theorem~\ref{thm:T3} is therefore excluding a class that is approached from
every direction but one, not quantifying over an empty set.
\end{remark}

\begin{remark}[a negative result worth recording]\label{rem:decreasing}
\Scomp{the tight instance of Remark~\ref{rem:T3num}}
The natural entry-wise bound $D_i \le d_i - (i-1)$ --- which, if true, would close the
whole problem by a one-line summation --- already fails \emph{inside} the hard core: for
$\dg = (3,3,3,3,3,2,1)$ one has $D = (3,3,2)$, so $d_2 - D_2 = 0$. Moreover the sum that a
proof of this shape would need is tight in that same instance. Any argument resting only on
entry-wise decrement estimates is therefore dead, which is why Lemma~\ref{lem:DICH}
proceeds instead by bounding head values through survivor decay.
\end{remark}

\section{A \texorpdfstring{$\tau$}{tau}-uniform budget}\label{sec:budget}

From here on the parameter that organises everything is $L = |\Blo|$, the number of
vertices of $B$ of degree at most $\tau$. Theorem~\ref{thm:K} says the hard core has
$L \ne 0$; this section proves $L \ge 2$ and Section~\ref{sec:fan} raises it to $L \ge 3$.
The results of this section carry registry rows R-16, R-17 and R-18. Note that the
high/low split is by \emph{degree}, not by head status; Lemma~\ref{lem:HI} makes the two
agree on $\Bhi$ only, and it is exactly that asymmetry which lets Theorem~\ref{thm:LOW}
say something when $K \ne B$.

\subsection{The bound}

\begin{lemma}\label{lem:HI}
\Scert{row R-16, tier ``reductio only''}
Under the reductio hypothesis, $\Bhi$ is exactly the set of heads of original degree
$\ge \tau+1$; in particular there are $p$ such heads and $L$ heads of original degree
$\le \tau$. Write $I_{\mathrm{hi}}$, $I_{\mathrm{lo}}$ for the corresponding sets of head
\emph{positions}, so $|I_{\mathrm{hi}}| = p$ and $|I_{\mathrm{lo}}| = L$.
\end{lemma}

\begin{proof}
Every $a \in A$ has $N(a) \subseteq B$, so $\dg(a) \le |B| = \tau$; hence every vertex of
degree $\ge \tau+1$ lies in $B$, and by Lemma~\ref{lem:S} every such vertex is a head.
\end{proof}

\begin{theorem}[the budget inequality]\label{thm:LOW}
\Scert{row R-16, tier ``reductio only''}
Under the reductio hypothesis $\res(G) = \alpha(G)$, with no hypothesis on $K$:
\begin{align}
\sum_{b \in \Blo} \dg(b) \;+\; \nu &\;\le\; L(\tau+1); \label{eq:LOW1}\\
\sum_{b \in \Blo} \degA(b) \;+\; e(\Blo) \;+\; \mbar &\;\le\; \tfrac{1}{2}L(L+3);
\label{eq:LOW2}\\
\sum_{b \in \Blo} \degA(b) \;+\; \mbar &\;\le\; 2L + \nu(\Blo). \label{eq:LOW3}
\end{align}
\end{theorem}

\begin{proof}
By Lemma~\ref{lem:hh}(2), $\sum_{i=1}^{\tau} D_i = m$. Split the head positions by
Lemma~\ref{lem:HI}. For $i \in I_{\mathrm{hi}}$, Lemma~\ref{lem:DICH}(b) gives
$D_i = g_i - (i-1)$ exactly, where $g_i$ is the original degree; for
$i \in I_{\mathrm{lo}}$, Lemma~\ref{lem:DICH}(c) gives $D_i \le \tau - i + 2$. Hence
\[
m \;\le\; \sum_{b \in \Bhi} \dg(b) \;-\; \sum_{i \in I_{\mathrm{hi}}}(i-1)
\;+\; \sum_{i \in I_{\mathrm{lo}}}(\tau - i + 2).
\]
On the other hand $m = \sum_{b \in B}\dg(b) - e_B$. Subtracting,
\[
\sum_{b \in \Blo}\dg(b) - e_B \;\le\;
- \sum_{i \in I_{\mathrm{hi}}}(i-1) + \sum_{i \in I_{\mathrm{lo}}}(\tau-i+2).
\]
The right-hand side is independent of which positions the low heads occupy: adding and
subtracting $\sum_{i \in I_{\mathrm{lo}}}(i-1)$,
\[
- \sum_{i \in I_{\mathrm{hi}}}(i-1) + \sum_{i \in I_{\mathrm{lo}}}(\tau-i+2)
= - \sum_{i=1}^{\tau}(i-1) + \sum_{i \in I_{\mathrm{lo}}}\bigl[(i-1)+(\tau-i+2)\bigr]
= - \binom{\tau}{2} + L(\tau+1),
\]
because the bracket collapses to the constant $\tau+1$. With $e_B = \binom{\tau}{2} - \nu$
this is \eqref{eq:LOW1}. For \eqref{eq:LOW2} write
$\sum_{\Blo}\dg(b) = \sum_{\Blo}\degA(b) + 2e(\Blo) + e(\Blo,\Bhi)$ and
$\nu = \nu(\Blo) + \nu(\Blo,\Bhi) + \mbar$ with
$\nu(\Blo) = \binom{L}{2} - e(\Blo)$ and $\nu(\Blo,\Bhi) = Lp - e(\Blo,\Bhi)$; the crossing
terms cancel and $L(\tau+1) - \binom{L}{2} - Lp = \tfrac12 L(L+3)$. Finally \eqref{eq:LOW3}
is \eqref{eq:LOW2} with $e(\Blo) = \binom{L}{2} - \nu(\Blo)$.
\end{proof}

Theorem~\ref{thm:K} is the case $L = 0$: then \eqref{eq:LOW1} reads $\nu \le 0$, i.e.\ $B$
is a clique. So Theorem~\ref{thm:LOW} subsumes it, and it never mentions $K$ at all.

\begin{theorem}[slack positivity]\label{thm:SL}
\Scert{row R-16, tier ``reductio only''}
Assume the reductio hypothesis and $\tau \ge 2$. If $L \ge 1$ then
\[
\mathrm{slack} \;:=\; L(\tau+1) - \Bigl( \sum_{b \in \Blo}\dg(b) + \nu \Bigr) \;\ge\; 1,
\]
i.e.\ \eqref{eq:LOW1} holds with one unit to spare; equivalently
$\sum_{b\in\Blo}\degA(b) + e(\Blo) + \mbar \le \tfrac12 L(L+3) - 1$.
\end{theorem}

\begin{proof}
From the proof of Theorem~\ref{thm:LOW},
$\mathrm{slack} = \sum_{i \in I_{\mathrm{lo}}}(\tau - i + 2 - D_i)$, each summand being
$\ge 0$ by Lemma~\ref{lem:DICH}(c). Suppose $\mathrm{slack} = 0$; then
\begin{equation}\label{eq:star}
\text{every low head satisfies } D_i = \tau - i + 2 \text{ exactly.}
\tag{$*$}
\end{equation}
Let $i_0 := \max I_{\mathrm{lo}}$ and let $x$ be the head deleted at step $i_0$, of original
degree $g = \dg(x) \le \tau$. Write $v_k$ for $x$'s value at the start of step $k$.

\emph{Step 1: $i_0 \ge 2$.} If $i_0 = 1$ then $I_{\mathrm{lo}} = \{1\}$, so every head at a
position $2,\dots,\tau$ is high, of original degree $\ge \tau+1$; but
$D_1 = \Delta = g \le \tau$ by Lemma~\ref{lem:hh}(2), contradicting that $\Delta$ is the
maximum degree. (If $L = \tau$ there are no high heads and $i_0 = \tau \ge 2$ directly.)

\emph{Step 2: an escape step exists, and it sits exactly at the decay threshold.} Let
$\mathcal{E} := \{k < i_0 : x \notin \mathrm{block}_k\}$. Since
$v_{i_0} = g - (i_0 - 1 - |\mathcal{E}|)$ and $v_{i_0} = D_{i_0} = \tau - i_0 + 2$ by
\eqref{eq:star}, we get $|\mathcal{E}| = \tau + 1 - g \ge 1$. Put $j := \max \mathcal{E}$.
From step $j+1$ to step $i_0-1$ the entry $x$ is decremented every time, so
$v_{i_0} = v_j - (i_0-1-j)$, whence $v_j = \tau - j + 1$.

\emph{Step 3: what $\mathrm{block}_j$ can contain.} $\mathrm{block}_j$ is the prefix of the
$D_j$ largest non-head entries in non-increasing order, and $x$ is an entry of step $j$
lying outside it; hence every entry of $\mathrm{block}_j$ has value $\ge v_j = \tau-j+1$. By
Lemma~\ref{lem:F3prime} every survivor has value $\le \tau-j+1$ at the start of step $j$.
Therefore each survivor in $\mathrm{block}_j$ has value exactly $\tau-j+1$ and so --- needing
$\tau-j+1$ further decrements in the $\tau-j+1$ remaining steps --- lies in
$\mathrm{block}_k$ for every $k \in [j,\tau]$. Write $s_j$ for the number of survivors in
$\mathrm{block}_j$.

\emph{Step 4: the head part of $\mathrm{block}_j$.} Every entry of $\mathrm{block}_j$ is
either a survivor or a head deleted later. Let $\ell_j := \#\{i \in I_{\mathrm{lo}} : i > j\}$
and note $x$ is one of these and $x \notin \mathrm{block}_j$, so the low later heads
contribute $\lambda_j \le \ell_j - 1$. Every high later head lies in every earlier block by
Lemma~\ref{lem:DICH}(b), and there are $(\tau-j)-\ell_j$ of them. Hence
$D_j = |\mathrm{block}_j| = (\tau-j-\ell_j) + \lambda_j + s_j$.

\emph{Step 5: $s_j \le 2$.} All heads after position $i_0$ are high, so by
Lemma~\ref{lem:DICH}(b) all $\tau - i_0$ of them lie in $\mathrm{block}_{i_0}$; by Step 3
the $s_j$ survivors lie in $\mathrm{block}_{i_0}$ too, and these sets are disjoint. Hence
$D_{i_0} \ge s_j + (\tau-i_0)$, and $D_{i_0} = \tau-i_0+2$ gives $s_j \le 2$.

\emph{Step 6: contradiction.} Steps 4 and 5 give
$D_j \le (\tau-j-\ell_j) + (\ell_j-1) + 2 = \tau-j+1$. But position $j < i_0$ is a head
position, and either it is high, so $D_j = g_j - (j-1) \ge (\tau+1)-(j-1) = \tau-j+2$ by
Lemma~\ref{lem:DICH}(b), or it is low, so $D_j = \tau-j+2$ by \eqref{eq:star}. Either way
$D_j \ge \tau-j+2 > \tau-j+1 \ge D_j$.
\end{proof}

\begin{remark}[sharpness]\label{rem:SLsharp}
\Scomp{$19\,710$ trajectories in all, of which $2\,740$ are on the exhaustive atlas
$n \le 7$ and the remainder on $9\,000$ random graphs $n = 8,\dots,13$ plus $4\,000$
near-split graphs, five trajectories each}
Theorem~\ref{thm:SL} is \emph{not} the statement ``the last low head satisfies
$D_{i_0} \le \tau - i_0 + 1$'': that stronger claim is false, with $1\,925$ counterexamples
in the corpus. What Step 6 extracts is that the escape step $j$ must itself be a low head
position carrying strict inequality, so equality cannot hold at all low heads
simultaneously --- which is $\mathrm{slack} \ge 1$ and no more. Every intermediate
assertion of the proof was checked separately, not just the conclusion, with $0$ failures,
and the histogram of $s_j$ at equality heads is $\{2 : 2300\}$: Step 5's bound is attained
every time it is used. In a separate corpus of $6\,403$ graphs and $34\,040$ reductio
trajectory checks, the minimum slack of \eqref{eq:LOW1} per $L$ is
$\{0:0,\,1:1,\,2:1,\,3:2,\,4:2,\,5:3,\dots\}$: tight exactly at $L = 0$, which is forced by
Theorem~\ref{thm:K}, and never $0$ for $L \ge 1$.
\end{remark}

\begin{corollary}[the hard core has $L \ge 2$]\label{cor:SLHC}
\Scert{row R-18}
In the hard core, at least two vertices of $B$ have degree $\le \tau$.
\end{corollary}

\begin{proof}
$L = 0$ is impossible: Theorem~\ref{thm:K} would make $B$ a clique, contradicting
Observation~\ref{obs:R1}. $L = 1$ is impossible by Corollary~\ref{cor:L1short} below (or,
by the longer route of Proposition~\ref{prop:L1} and Corollary~\ref{cor:L1prime}, because
the slack would be exactly $0$, contradicting Theorem~\ref{thm:SL}). The numerical
hypothesis $\tau \ge 2$ of Theorem~\ref{thm:SL} is supplied by $\diam = 4$ alone: a
connected graph with $\tau \le 1$ has a vertex cover of size $\le 1$, hence is edgeless or a
star, hence has $\diam \le 2$.
\end{proof}

The last sentence of that proof is worth a word. In three successive states of this text the
parenthetical there cited $\tau \ge 3$, then $\tau \ge 4$, then Theorem~\ref{thm:T3} --- each
time a true statement that the proof does not consume, and each time making a self-contained
corollary depend on the heaviest theorem in the line, hence unreviewable by any judge not
handed that theorem. A review round objected, and the elementary local replacement above is
the repair. The species --- \emph{a true import nobody needs} --- recurred six times in the
campaign's source report and is the single most common defect its review gate caught.

\subsection{The master budget}

\begin{theorem}[master budget]\label{thm:MB}
\Scert{row R-17}
In the hard core, for $L \ge 1$,
\[
|\Bloplus| + c + \mbar \;\le\; L + \nu(\Blo) - 1,
\qquad\text{where } c := |\Blo \cap (T_1 \cup T_2)| .
\]
\end{theorem}

\begin{proof}
Lemma~\ref{lem:pair} gives $\degA(b) \ge 2$ for every $b \in \Bloplus$, and
$\degA(b) \ge 1$ for every $b \in \Blo$ since $A$ is a maximum independent set.
Lemma~\ref{lem:Cstar} gives $\degA(b) \ge 3$ for $b \in T_1 \cup T_2$, and every such $b$
\emph{that is low} lies in $\Bloplus$ --- it has the whole of the other type among its
non-neighbours --- so $\Blo \cap (T_1 \cup T_2) \subseteq \Bloplus$ and $c \le |\Bloplus|$.
Hence
\[
\sum_{b \in \Blo}\degA(b) \;\ge\; (L - |\Bloplus|)\cdot 1 + (|\Bloplus| - c)\cdot 2
+ c \cdot 3 \;=\; L + |\Bloplus| + c .
\]
Substituting into the \eqref{eq:LOW3} form of Theorem~\ref{thm:SL},
$\sum_{b\in\Blo}\degA(b) + \mbar \le 2L + \nu(\Blo) - 1$, gives the claim.
\end{proof}

The restriction ``that is low'' in the third sentence is not decoration: the corresponding
sentence without it is false for a \emph{high} $b \in T_1 \cup T_2$, since
$\Bloplus \subseteq \Blo$, and Proposition~\ref{prop:L2}(d) below is its own witness. That
was a review finding; the count is unaffected, since it quantifies only over $b \in \Blo$.

\begin{corollary}\label{cor:L1short}
\Scert{row R-17}
The hard core has no instance with $L = 1$.
\end{corollary}

\begin{proof}
$L = 1$ gives $\nu(\Blo) = 0$, so Theorem~\ref{thm:MB} reads $|\Bloplus| + c + \mbar \le 0$:
the unique low vertex $b_0$ is $B$-universal, $c = 0$ and $\mbar = 0$. $B$-universality of
$b_0$ means no non-edge of $B$ meets $b_0$, so $\nu = \mbar = 0$ and $B$ is a clique,
contradicting Observation~\ref{obs:R1}.
\end{proof}

\begin{corollary}[a floor on $L$]\label{cor:MB1}
\Scert{row R-17}
In the hard core, if every low vertex is $B$-universal then $\nu = \mbar \le L - 1$, so
$L \ge \nu + 1 \ge 2$.
\end{corollary}

\begin{proof}
$B$-universality of all of $\Blo$ gives $\nu(\Blo) = \nu(\Blo,\Bhi) = 0$, hence
$\nu = \mbar$, and $c = 0$ since a $B$-universal vertex has no non-neighbour and so cannot
lie in $T_1 \cup T_2$. Theorem~\ref{thm:MB} then reads $\mbar \le L-1$, and $\nu \ge 1$ by
Observation~\ref{obs:R1}.
\end{proof}

\subsection{The two rigid low configurations}

The following two statements were the campaign's original route to $L \ge 2$; they are
superseded in part by Corollary~\ref{cor:L1short} and Theorem~\ref{thm:RIG} but are
certified in their own right, and Proposition~\ref{prop:L2} is what identifies the
configuration that Section~\ref{sec:fan} then eliminates.

\begin{proposition}[$L = 1$ is rigid]\label{prop:L1}
\Scert{row R-17}
In the hard core with $L = 1$, write $\Blo = \{b_0\}$. Then $b_0$ is adjacent to every
other vertex of $B$, $\degA(b_0) = 1$, $\dg(b_0) = \tau$ and $\nu = \mbar = 1$; the unique
non-adjacent pair of $B$ is a pair $u,v \in \Bhi$ with $T_1 = \{u\}$, $T_2 = \{v\}$ and
$\degA(u), \degA(v) \ge 3$; and the unique $A$-neighbour $a_0$ of $b_0$ is adjacent to all
of $B$.
\end{proposition}

\begin{proof}
$L = 1$ gives $\nu(\Blo) = 0$, so \eqref{eq:LOW3} of Theorem~\ref{thm:SL} reads
$\degA(b_0) + \mbar \le 2$. If $b_0 \in T_1 \cup T_2$ then Lemma~\ref{lem:Cstar} gives
$\degA(b_0) \ge 3$, impossible; so $T_1, T_2 \subseteq \Bhi$. If $b_0$ had a non-neighbour
in $B$, Lemma~\ref{lem:pair} would give $\degA(b_0) \ge 2$, hence $\mbar = 0$, i.e.\ $\Bhi$
is a clique; but $T_1, T_2$ are disjoint non-empty subsets of $\Bhi$, so some $u \in T_1$
and $v \in T_2$ are adjacent, contradicting Fact~\ref{fact:Fb}. So $b_0$ is adjacent to all
of $B \setminus \{b_0\}$, whence $\nu = \mbar$; by Observation~\ref{obs:R1} $\nu \ge 1$, so
$\degA(b_0) + \mbar \le 2$ forces $\degA(b_0) = 1$ and $\nu = \mbar = 1$, and then
$\dg(b_0) = (\tau-1)+1 = \tau$. Since the only non-adjacent pair of $B$ is that single pair
$uv \subseteq \Bhi$, and all cross pairs of $T_1 \times T_2$ must be non-adjacent,
$|T_1| = |T_2| = 1$ and $\{T_1,T_2\} = \{\{u\},\{v\}\}$; Lemma~\ref{lem:Cstar} gives
$\degA(u),\degA(v) \ge 3$. Finally, for each $w \in B \setminus \{b_0\}$ the pair
$(b_0,w)$ is adjacent, so Lemma~\ref{lem:pair} gives a common $A$-neighbour, which must be
$a_0$.
\end{proof}

\begin{corollary}\label{cor:L1prime}
\Scert{row R-17}
In the hard core with $L = 1$ the slack of \eqref{eq:LOW1} is exactly $0$:
$\sum_{b \in \Blo}\dg(b) + \nu = \tau + 1 = L(\tau+1)$.
\end{corollary}

\begin{proof}
Immediate from Proposition~\ref{prop:L1}: $\dg(b_0) = \tau$ and $\nu = 1$.
\end{proof}

\begin{proposition}[$L = 2$ is rigid]\label{prop:L2}
\Scert{row R-17}
Suppose the hard core has an instance with $L = 2$, $\Blo = \{b_1,b_2\}$. Then:
\begin{enumerate}[label=\textup{(\alph*)}]
\item $b_1 \sim b_2$;
\item both are $B$-universal;
\item $\degA(b_1) = \degA(b_2) = 1$ and their unique $A$-neighbours coincide in a single
vertex $a_0$ adjacent to all of $B$;
\item $\nu = \mbar = 1$, i.e.\ $G[B] = K_\tau$ minus exactly one edge $uv$ with
$u,v \in \Bhi$, and $T_1 = \{u\}$, $T_2 = \{v\}$, $\degA(u),\degA(v) \ge 3$;
\item the slack of \eqref{eq:LOW1} is exactly $1$, so Theorem~\ref{thm:SL} is tight here.
\end{enumerate}
\end{proposition}

\begin{proof}
Theorem~\ref{thm:MB} reads $|\Bloplus| + c + \mbar \le 1 + \nu(\Blo)$.

(a) If $b_1 \not\sim b_2$ then $\nu(\Blo) = 1$ and both lie in $\Bloplus$, so
$2 + c + \mbar \le 2$, giving $c = \mbar = 0$. Then $\Bhi$ is a clique and
$T_1, T_2 \subseteq \Bhi$ are disjoint and non-empty, so some $u \in T_1$ is adjacent to
some $v \in T_2$, contradicting Fact~\ref{fact:Fb}. Hence $b_1 \sim b_2$ and
$\nu(\Blo) = 0$, so $|\Bloplus| + c + \mbar \le 1$.

(b) If some $b_i \in \Bloplus$ then $c = \mbar = 0$ and the same contradiction recurs.
Hence $\Bloplus = \emptyset$ and $c = 0$, leaving $\mbar \le 1$.

(d) Both $b_i$ being $B$-universal, every non-edge of $B$ lies inside $\Bhi$, so
$\nu = \mbar \le 1$; Observation~\ref{obs:R1} gives $\nu \ge 1$, so $\nu = \mbar = 1$. The
unique non-adjacent pair of $B$ is $uv \subseteq \Bhi$; since all cross pairs of
$T_1 \times T_2$ are non-adjacent, $T_1 = \{u\}$ and $T_2 = \{v\}$, and
Lemma~\ref{lem:Cstar} gives $\degA(u),\degA(v) \ge 3$.

(c) At $L = 2$, $|\Bloplus| = c = 0$, $\mbar = 1$, $\nu(\Blo) = 0$, both bounds combined in
the proof of Theorem~\ref{thm:MB} are simultaneously tight: the degree count gives the lower
bound $\sum_{b\in\Blo}\degA(b) \ge L + |\Bloplus| + c = 2$, while \eqref{eq:LOW3} of
Theorem~\ref{thm:SL} gives the upper bound
$\sum_{b\in\Blo}\degA(b) \le 2L + \nu(\Blo) - 1 - \mbar = 2$. The two meet, so
$\degA(b_1) = \degA(b_2) = 1$. Since $b_1 \sim b_2$, Lemma~\ref{lem:pair} supplies a common
$A$-neighbour, necessarily the unique $A$-neighbour of each; call it $a_0$. For every
$w \in B \setminus \{b_1\}$ the pair $(b_1,w)$ is adjacent, so Lemma~\ref{lem:pair} gives a
common $A$-neighbour, necessarily $a_0$; hence $a_0$ is adjacent to all of $B$.

(e) $\dg(b_i) = (\tau-1)+1 = \tau$, so
$\mathrm{slack} = 2(\tau+1) - (2\tau + \nu) = 1$.
\end{proof}

Writing the two bounds of (c) out as a squeeze, rather than saying only ``re-run the degree
count'', was a review repair: a tightness argument needs the inequality chain to be tight at
every link, and naming one of the two links is not enough.

\section{The fan family}\label{sec:fan}

Proposition~\ref{prop:L2} says that the $L = 2$ layer of the hard core, if non-empty, is a
single rigid configuration, on which Theorem~\ref{thm:SL} is tight and therefore cannot
help. This section eliminates that configuration --- and its whole family --- by a
Havel--Hakimi argument.

\begin{definition}[the fan configurations]\label{def:fan}
For $\tau \ge 2$, $L \ge 1$ and $\nu \ge 1$, say that $G$ carries the configuration
$\GFan(\tau,L,\nu)$ if: $|B| = \tau$; the $L$ vertices of $\Blo$ are all $B$-universal, each
with $\degA = 1$, and their common $A$-neighbour is a single vertex $a_0$ adjacent to all
of $B$; all $\nu$ non-edges of $B$ lie inside $\Bhi$; and every $a \in A' := A \setminus
\{a_0\}$ has $N(a) \subseteq \Bhi$. Write $\Fan(\tau,L) := \GFan(\tau,L,1)$, so that
$\Fan(\tau,L)$ is the configuration of Proposition~\ref{prop:L2} at general $L$:
$G[B] = K_\tau$ minus one edge $uv$ with $u,v \in \Bhi$.
Throughout put $C := \Blo \cup \{a_0\}$ and $p := \tau - L = |\Bhi|$.
\end{definition}

In $\Fan(\tau,L)$ every $c \in C$ has $\dg(c) = \tau$ exactly ($\tau-1$ inside $B$ plus
$a_0$ for $c \in \Blo$; and $a_0$ has exactly its $\tau$ neighbours in $B$); every
$w \in \Bhi$ has $\dg(w) \ge \tau+1$; and every $a \in A'$ has $\dg(a) \le p$. Note
$u,v \in \Bhi$ forces $p \ge 2$, i.e.\ $L \le \tau-2$.

We call steps $1,\dots,p$ the \emph{high phase}, and say a vertex \emph{escapes}
$\mathrm{block}_t$ if it is not decremented at step $t$; $E$ denotes the total number of
$(\text{vertex},\text{step})$ escapes of $C$-vertices during the high phase.

\begin{lemma}[high phase first]\label{lem:FAN1}
\Scert{row R-9}
In $\Fan(\tau,L)$ under the reductio hypothesis, the first $p$ heads may be taken to be
exactly the $p$ vertices of $\Bhi$.
\end{lemma}

\begin{proof}
Swapping the roles of two entries of equal current value --- one taken as head, the other as
a block member --- leaves the value multiset trajectory unchanged, hence leaves $s$
unchanged; so the reductio hypothesis is tie-break invariant and we may fix the tie-break.
Induct on $j \le p$: if all heads before $j$ are high, then $p - (j-1) \ge 1$ high vertices
remain, each in every earlier block by Lemma~\ref{lem:DICH}(b), hence of value
$\dg - (j-1) \ge \tau-j+2$, so $D_j \ge \tau-j+2$. A low head at $j$ has
$D_j \le \tau-j+2$ by Lemma~\ref{lem:DICH}(c); equality is a tie with a remaining high
vertex, and the fixed tie-break takes the high vertex.
\end{proof}

\begin{lemma}[the endgame is unconditional]\label{lem:FAN3}
\Scert{row R-9}
If at the start of some step the remaining multiset is $[L]^{L+1},1,1$ with $L \ge 2$, then
the process needs exactly $L+1$ further steps.
\end{lemma}

\begin{proof}
Head values are non-increasing. At the first of these steps the maximum is $L$ and the $L$
other $L$-entries are the top $L$ non-head entries, so the block is exactly those, they
become $L-1$, and the two $1$s survive untouched: the list becomes $[L-1]^{L},1,1$.
Inductively after $j$ steps the list is $[L-j]^{L+1-j},1,1$, so after $L-1$ steps it is
$[1]^{2},1,1 = [1,1,1,1]$: the head is $1$, $D = 1$, one entry is zeroed and $[0,1,1]$
remains. One further step clears it. That is $L+1$ steps in total.
\end{proof}

\begin{lemma}[the residue mass]\label{lem:FAN4}
\Scert{row R-9}
In $\Fan(\tau,L)$ under the reductio hypothesis, assume $E = 0$. Then the $A'$ entries at
the start of step $p+1$ sum to exactly $2$.
\end{lemma}

\begin{proof}
By Lemma~\ref{lem:FAN1} the heads of steps $1,\dots,p$ are $\Bhi$, so by
Lemma~\ref{lem:DICH}(b),
$\sum_{j \le p} D_j = \sum_{x \in \Bhi}\dg(x) - \binom{p}{2}$. Now
$\sum_{x \in \Bhi}\degB(x) = 2(\tau-2) + (p-2)(\tau-1) = p\tau - p - 2$, since $u,v$ miss
the edge $uv$ and the other $p-2$ vertices of $\Bhi$ do not, and
$\sum_{x \in \Bhi}\degA(x) = p + R$ where $R := \sum_{a \in A'}\dg(a)$ and the $p$ counts
$a_0$'s edges to $\Bhi$; so $\sum_{\Bhi}\dg = p\tau - 2 + R$. Counting the same decrements
by recipient: a high head deleted at position $i \le p$ receives $i-1$ of them, total
$\binom{p}{2}$; each of the $L+1$ vertices of $C$ receives $p$ by hypothesis, total
$p(L+1)$; the rest go to $A'$. Hence
$\mathrm{dec}_{A'} = (p\tau - 2 + R - \binom{p}{2}) - \binom{p}{2} - p(L+1) = R - 2$, and
$A'$ starts at total $R$, leaving total $2$.
\end{proof}

So the residual $A'$ partition is either $1+1$ or a single $2$, and nothing else is
arithmetically possible. Which of the two occurs decides the whole family:

\begin{observation}\label{obs:FAN5}
\Scert{row R-14}
For every $L \ge 2$, the multiset $[L]^{L+1},1,1$ clears in $L+1$ Havel--Hakimi steps, while
$[L]^{L+1},2$ clears in exactly $L$.
\end{observation}

\begin{proof}
The first half is Lemma~\ref{lem:FAN3}. The second is Lemma~\ref{lem:TAIL} below at
$\lambda = [2]$: there $\lambda_1 = 2$ and $s_0(\lambda) = \mathrm{steps}([2]^3 \cup [2]) =
2 = \lambda_1$, so the total is $(L-2) + 2 = L$ for every $L \ge 2$.
\end{proof}

Since the reductio hypothesis requires exactly $L$ further steps after step $p$, an instance
whose $A'$ residue were a single $2$ would be a genuine survivor. That is the single
load-bearing point of the whole elimination, at every $L \ge 2$ --- not a technicality of
$L = 2$, as an earlier reading had it.

\begin{lemma}[escape budget]\label{lem:FAN7}
\Scert{row R-9}
In $\Fan(\tau,L)$ under the reductio hypothesis, $E \le 1$.
\end{lemma}

\begin{proof}
Repeating the count of Lemma~\ref{lem:FAN4} with $p(L+1) - E$ decrements delivered to $C$
gives $\mathrm{dec}_{A'} = R - 2 + E$, so the $A'$ residue is $2 - E \ge 0$, whence
$E \le 2$. If $E = 2$ the $A'$ residue is $0$ and the $C$-part at the start of step $p+1$ is
$[L+2, L^L]$ or $[L+1,L+1,L^{L-1}]$; in both, the head exceeds the number of positive
remaining entries ($L+2 > L$, resp.\ $L+1 > L$), so the next step would drive a zero entry
negative, which is impossible for the degree sequence of a graph.
\end{proof}

\begin{lemma}[no escape at all]\label{lem:FAN8}
\Scert{row R-9}
In $\Fan(\tau,L)$ with $L \ge 2$, under the reductio hypothesis, $E = 0$.
\end{lemma}

\begin{proof}
Suppose $E \ge 1$ and let $t \le p$ be a step at which $e_t \ge 1$ vertices of $C$ escape;
let $c$ be one of them. By Lemma~\ref{lem:FAN1} the head of step $t$ is a $\Bhi$ vertex, so
$D_t = g_t - (t-1) \ge \tau-t+2$; by Lemma~\ref{lem:DICH}(b) the $p-t$ remaining $\Bhi$
vertices lie in $\mathrm{block}_t$, and so does each of the $(L+1)-e_t$ non-escaping
$C$-vertices. Hence
$|\mathrm{block}_t \cap A'| = D_t - (p-t) - (L+1) + e_t \ge (\tau-t+2) - (p-t) - (L+1) + 1
= 2$, so $\mathrm{block}_t$ contains some $x \in A'$. Because $\mathrm{block}_t$ is a prefix
of the sorted list and $c \notin \mathrm{block}_t$, we get
$v_t(x) \ge v_t(c) \ge \tau-t+1$. By the count in Lemma~\ref{lem:FAN7} the $A'$ entries sum
to $2 - E \le 1$ at the start of step $p+1$, so $v_{p+1}(x) \le 1$; and $x$ is not deleted
in steps $t,\dots,p$ (all those heads are $\Bhi$), so it loses at most one per step and
$v_{p+1}(x) \ge v_t(x) - (p-t+1)$. Therefore $\tau-t+1 \le 1 + (p-t+1)$, i.e.\
$\tau \le p+1 = \tau-L+1$, i.e.\ $L \le 1$.
\end{proof}

\begin{lemma}[backward induction]\label{lem:FAN6}
\Scert{row R-9}
In $\Fan(\tau,L)$ under the reductio hypothesis, assume $E = 0$. Then the $A'$ residue at
the start of step $p+1$ is $1+1$, not a single $2$.
\end{lemma}

\begin{proof}
Write $M_j$ for the $A'$ value multiset at the start of step $j$. Since all of $\Bhi$
(Lemma~\ref{lem:DICH}(b)) and all of $C$ (hypothesis) lie in $\mathrm{block}_j$, the number
of $A'$ entries in $\mathrm{block}_j$ is $a_j = D_j - (p-j) - (L+1) = g_j - \tau \ge 1$, and
--- the block being a prefix of the sorted list --- they are the $a_j$ largest entries of
$M_j$. In particular a maximum entry of $M_j$ is decremented at every step. By
Lemma~\ref{lem:FAN4} the total of $M_{p+1}$ is $2$, so it is $1+1$ or a single $2$; suppose
it is a single $2$. Say $M_{j+1}$ has a unique maximum $v$ and does not contain the value
$v-1$. At step $j$ let $M_j$ have maximum $w$ of multiplicity $\mu$. If $a_j < \mu$ the
maximum does not drop, $w = v$, and the $a_j \ge 1$ decremented copies land at $v-1$, which
$M_{j+1}$ does not contain --- a contradiction. So $a_j \ge \mu$: all maxima drop,
$w = v+1$, and $M_{j+1}$ contains $\mu$ copies of $v$, forcing $\mu = 1$. Every other entry
of $M_j$ is either in the block (its $M_{j+1}$ value plus $1$) or out of it (unchanged), so
all are at most $\max(\text{other entries of } M_{j+1}) + 1$. Hence the hypothesis
propagates backwards: $M_{p+1} = [2]$ has unique maximum $2$ and no $1$, so $M_p = [3,1^b]$,
$M_{p-1}$ has unique maximum $4$ with the others $\le 2$, and inductively $M_{p-t}$ has
maximum $3+t$ with all other entries $\le t+1$ --- in particular $v-1$ is never present, so
the induction never stops. At $t = p-1$ this gives $\max(M_1) = p+2$. But $M_1$ is the
multiset of $A'$ \emph{degrees}, and every $a \in A'$ has $N(a) \subseteq \Bhi$, so
$\max(M_1) \le p$.
\end{proof}

The induction does not apply to $1+1$: its maximum is not unique --- the value $1$ with
multiplicity $2$ --- and the first step needs a unique maximum. That is exactly why the
lemma kills only the single-$2$ case, and it is not a gap. (An earlier text gave a different
and false reason, namely that the value $v-1 = 0$ is present among the zero entries; that
fails whenever $|A'| = 2$, which is realisable at every $L \ge 2$.)

\begin{theorem}[the fan family is eliminated]\label{thm:FAN}
\Scert{row R-10}
For every $\tau$ and every $L \ge 2$ there is no graph carrying the configuration
$\Fan(\tau,L)$ and satisfying $\res(G) = \alpha(G)$. Consequently, by
Fact~\ref{fact:fms}, every $\Fan(\tau,L\ge2)$ graph has $\res \le \alpha - 1$.
\end{theorem}

\begin{proof}
Assume $\res = \alpha$, i.e.\ $s = \tau$. By Lemma~\ref{lem:FAN1} the heads of steps
$1,\dots,p$ are exactly $\Bhi$. By Lemma~\ref{lem:FAN8}, $E = 0$, so by
Lemma~\ref{lem:FAN4} the $A'$ entries sum to exactly $2$ at the start of step $p+1$, and by
Lemma~\ref{lem:FAN6} they are $1+1$; the $C$-entries are then all at $\tau - p = L$. So the
list at the start of step $p+1$ is $[L]^{L+1},1,1$ (plus zeros), and
Lemma~\ref{lem:FAN3} says $L+1$ further steps are needed, giving $s = \tau+1$ --- a
contradiction.
\end{proof}

The conclusion is stated as $\res \le \alpha-1$ and not as $s = \tau+1$. The stronger reading
is available only \emph{inside} the refuted hypothesis and may not be carried out of it;
an earlier version of this theorem carried it out, which a review round caught. What is
true, and measured, is stated separately:

\begin{observation}\label{obs:FANE}
\Scomp{$1\,098\,141$ strict and $1\,700\,094$ superset $\Fan(\tau,L\ge2)$ degree sequences
in a box $\tau \le 10$, $R \le 20$, $|A'| \le 14$}
In every one of those instances $\res = \alpha - 1$ and $s = \tau+1$. This is a
\emph{lead}, not a theorem: at $s = \tau+1$ the reductio-only toolkit
(Lemmas~\ref{lem:S} and~\ref{lem:DICH}(b)) is unavailable --- witness $\Fan(4,2)$ with
$A' = [2,1,1]$, $\dg = [5,5,4,4,4,2,1,1]$, where every $\Bhi$ degree is
$\tau+1 = 5 = s$. Separately, over $508\,239$ labelled high-phase runs under adversarial
tie-breaks, $E = 0$ and the residue partition is $(1,1)$ in all $508\,239$; and over
$450\,303$ superset sequences the multiset at the start of step $p+1$ is exactly
$[L]^{L+1}$ with $A'$ residue $(1,1)$ every time, the partition $(2)$ never occurring.
These runs are mechanism checks off the reductio hypothesis in the sense of
Section~\ref{sec:conventions}; the lemmas rest on their proofs alone.
\end{observation}

\begin{corollary}[the hard core has $L \ge 3$]\label{cor:FANHC}
\Scert{row R-19}
In the hard core the number of $B$-vertices of degree $\le \tau$ satisfies $L \ge 3$;
together with Theorem~\ref{thm:T3}, the residual hard core has $\tau \ge 4$ and $L \ge 3$.
\end{corollary}

\begin{proof}
$L = 0$ is Theorem~\ref{thm:K} with Observation~\ref{obs:R1}; $L = 1$ is
Corollary~\ref{cor:L1short}. For $L = 2$: by Proposition~\ref{prop:L2} every hard-core
instance with $L = 2$ is rigidly $\Fan(\tau,2)$ --- clauses (a)--(d) supply every clause of
Definition~\ref{def:fan}, the clause on $A'$ coming from $\degA(b_i) = 1$ --- and
Theorem~\ref{thm:FAN} says no such graph satisfies $\res = \alpha$.
\end{proof}

The asymmetry here is deliberate. Theorem~\ref{thm:FAN} kills the whole family
$\Fan(\tau,L)$ for every $L \ge 2$, but it removes only the $L = 2$ layer of the hard core,
because rigidity is available at $L = 2$ only. For $L \ge 3$ a hard-core instance need not
be a fan; Section~\ref{sec:rig} supplies the rigidity that is available.

\subsection{A tail lemma}

\begin{lemma}[tail]\label{lem:TAIL}
\Scert{row R-15}
Let the multiset at the start of some step be $[L]^{L+1} \cup \lambda$ with $\lambda$ a
non-empty partition, $\lambda_1 := \max \lambda$. If $L \ge \lambda_1$ then the number of
further Havel--Hakimi steps is $(L - \lambda_1) + s_0(\lambda)$, where
$s_0(\lambda) := \mathrm{steps}([\lambda_1]^{\lambda_1+1} \cup \lambda)$. Consequently,
still under $L \ge \lambda_1$, ``clears in exactly $L$ steps'' is equivalent to
$s_0(\lambda) = \lambda_1$, a condition on $\lambda$ alone.
\end{lemma}

\begin{proof}
If $L = \lambda_1$ the list already is $[\lambda_1]^{\lambda_1+1} \cup \lambda$ and the
count reads $0 + s_0(\lambda)$; so assume $L > \lambda_1$. Suppose at some stage the list is
$[t]^{t+1} \cup \lambda$ with $t > \lambda_1$, which holds at $t = L$. The head is $t$; the
remaining entries of value $t$ are exactly the other $t$ copies, and every $\lambda$-entry
is $< t$; so the block --- the top $t$ non-head entries --- is exactly those $t$ copies,
they become $t-1$, and $\lambda$ is untouched. Hence
$[t]^{t+1} \cup \lambda \to [t-1]^{t} \cup \lambda$ in one step. Iterating from $t = L$
down to $t = \lambda_1$ uses $L - \lambda_1$ steps and leaves
$[\lambda_1]^{\lambda_1+1} \cup \lambda$.
\end{proof}

Three hypotheses in that statement are there because review rounds put them there:
$\lambda \ne \emptyset$, because $\lambda_1$ is undefined otherwise; the case $L = \lambda_1$
treated separately, because the induction's premise $t > \lambda_1$ is false at $t = L$
there; and the standing range $L \ge \lambda_1$ attached to the ``$L$-free'' consequence,
which is false without it --- $\lambda = [2]$, $L = 0$ has $s_0 = 2 = \lambda_1$ yet takes
one step, and $\lambda = [3]$, $L = 1$ clears in $L = 1$ step although
$s_0([3]) = 4 \ne 3$. Every use of the lemma in this paper is in range.

\begin{remark}\label{rem:TAILnum}
\Scomp{$1\,817$ pairs $(\lambda,L)$ with $\lambda \vdash 2\nu$, $\nu \le 6$, and
$\lambda_1 \le L \le 15$}
The predicted count $(L-\lambda_1)+s_0(\lambda)$ was compared against direct simulation on
all of them, with $0$ mismatches. Observation~\ref{obs:FAN5}'s second half was
independently confirmed for $L = 2,\dots,9$ by direct simulation, and a reviewing session
supplied a second, independent proof of it.
\end{remark}

\section{Rigidity of the \texorpdfstring{$B$}{B}-universal layer}\label{sec:rig}

Corollary~\ref{cor:FANHC} leaves $L \ge 3$, where no analogue of
Proposition~\ref{prop:L2} was available. This section supplies one. It rests on a
one-line degree cap, and it is the point at which the residual configuration space acquires
a name.

\begin{lemma}[degree cap on low vertices]\label{lem:CAP}
\Scert{row R-20}
Let $A$ be a maximum independent set, $B = V \setminus A$, $\tau = |B|$, and let
$b \in \Blo$ have $n_b$ non-neighbours in $B$. Then $1 \le \degA(b) \le n_b + 1$.
\end{lemma}

\begin{proof}
$\degA(b) \ge 1$: otherwise $N(b) \subseteq B$, so $A \cup \{b\}$ is independent,
contradicting maximality of $A$. For the upper bound, $\degB(b) = (\tau-1) - n_b$, and
$b \in \Blo$ means $\dg(b) \le \tau$, so
$\degA(b) = \dg(b) - \degB(b) \le \tau - (\tau-1-n_b) = n_b+1$.
\end{proof}

The hypotheses used are worth separating, and this is the one place in the campaign's
report that advertises minimality of hypotheses, so a review round tested it there. The
\emph{lower} bound uses only that $A$ is \emph{maximal} --- cardinality is never invoked ---
and the \emph{upper} bound uses neither maximality nor maximum cardinality, only the
definitions of $\Blo$, $\tau$ and $n_b$. So Lemma~\ref{lem:CAP} holds verbatim for any
maximal independent set. We state it with ``maximum'' because every consumer supplies one
anyway and because $\alpha = |A|$ is used in the same breath.

\begin{corollary}\label{cor:CAP1}
\Scert{row R-20}
A $B$-universal low vertex has $\degA(b) = 1$ exactly.
\end{corollary}

\begin{proof}
$n_b = 0$ in Lemma~\ref{lem:CAP}.
\end{proof}

\begin{theorem}[rigidity]\label{thm:RIG}
\Scert{row R-21}
Work in the hard-core frame and suppose $\Blo \ne \emptyset$ and every low vertex is
$B$-universal, i.e.\ $\Bloplus = \emptyset$. Then:
\begin{enumerate}[label=\textup{(\alph*)}]
\item $\degA(b) = 1$ for every $b \in \Blo$;
\item the unique $A$-neighbours of the low vertices all coincide in one vertex $a_0$;
\item $a_0$ is adjacent to all of $B$;
\item every non-edge of $B$ lies inside $\Bhi$, i.e.\ $\nu = \mbar$, and $\nu \ge 1$;
\item every $a \in A \setminus \{a_0\}$ has $N(a) \cap \Blo = \emptyset$.
\end{enumerate}
Consequently $G$ carries \emph{exactly} the configuration $\GFan(\tau,L,\nu)$ of
Definition~\ref{def:fan}, with $\nu = \mbar \ge 1$. In the hard core --- that is, adding the
reductio hypothesis --- one has in addition $\nu \le L-1$.
\end{theorem}

\begin{proof}
(a) is Corollary~\ref{cor:CAP1}. For (b), any two low vertices are adjacent by
$B$-universality, so Lemma~\ref{lem:pair} gives them a common $A$-neighbour, which by (a)
must be the unique $A$-neighbour of each; hence all coincide, in a vertex $a_0$. (For
$L = 1$ take $a_0$ to be the unique $A$-neighbour of the unique low vertex.) For (c), fix
$b \in \Blo$; for every $w \in B \setminus \{b\}$ we have $b \sim w$, so
Lemma~\ref{lem:pair} gives a common $A$-neighbour of $b$ and $w$, which by (a) is $a_0$;
hence $a_0 \sim w$ for all $w \in B \setminus \{b\}$, and $a_0 \sim b$ by definition. For
(d), a non-edge of $B$ incident to a low vertex would exhibit a non-neighbour of that vertex
in $B$, contradicting $B$-universality; so every non-edge lies inside $\Bhi$ and
$\nu = \mbar$, and $\nu \ge 1$ is Observation~\ref{obs:R1}. For (e), if $a \ne a_0$ had a
neighbour $b \in \Blo$ then $\degA(b) \ge 2$, contradicting (a).

The clauses of Definition~\ref{def:fan} are exactly (a)+(b)+(c), (d), (e), together with
``$\Bhi$ is high'', which is the definition of $\Bhi$. Conversely, $\GFan(\tau,L,\nu)$
\emph{is} the conjunction of exactly those clauses and asserts nothing else, so the match is
an equality and not merely an implication --- which is what ``exactly'' claims. The final
bound $\nu \le L-1$ is Corollary~\ref{cor:MB1} applied to this hypothesis, and that
corollary rests on Theorem~\ref{thm:MB} and hence Theorem~\ref{thm:SL}, which live in the
hard core.
\end{proof}

\begin{remark}[scope, stated exactly]\label{rem:RIGscope}
Clauses (a)--(c) and (e) need only that $A$ is a maximum independent set, plus
Lemma~\ref{lem:pair}; clause (d) needs additionally $\diam = 4$, through
Observation~\ref{obs:R1}. So (a)--(e) hold in the hard-core frame, with no reductio
hypothesis. Only the bound $\nu \le L-1$ needs the hard core, and that is not a formality:
see Remark~\ref{rem:RIGnum}. An earlier version of this paragraph glossed
Lemma~\ref{lem:pair} as though its only hypothesis were $f = \alpha+1$, and thereby claimed
that (a)--(c),(e) survive outside the frame. \emph{That gloss is not licensed.}
Lemma~\ref{lem:pair} as certified is a statement in the frame, and nothing in the campaign's
work establishes a version of it assuming $f = \alpha+1$ alone. The gloss was an affirmative
hypothesis-weakening claim about an import, which is exactly the class of claim the
campaign's scope discipline exists to police; a review round found it, and it is withdrawn
here. Theorem~\ref{thm:RIG}'s own statement and proof are unaffected.
\end{remark}

\begin{corollary}[the $B$-universal layer collapses to $\nu \ge 2$]\label{cor:RIG2}
\Scert{row R-21}
In the hard core with $L \ge 2$ and every low vertex $B$-universal,
Theorem~\ref{thm:RIG} gives $\GFan(\tau,L,\nu)$ with $1 \le \nu \le L-1$; the case
$\nu = 1$ is exactly $\Fan(\tau,L)$, which Theorem~\ref{thm:FAN} eliminates. Hence that
layer reduces to $\GFan(\tau,L,\nu)$ with $2 \le \nu \le L-1$ --- in particular it is empty
for $L = 2$, and at $L = 3$ it reduces to the single configuration $\GFan(\tau,3,2)$.
\end{corollary}

\begin{remark}\label{rem:RIGnum}
\Scomp{$2\,931$ pairs (graph, maximum independent set) over the exhaustive atlas
$n \le 7$, $4\,382$ over $1\,200$ random connected graphs $n = 8,\dots,10$, $1\,737$ over
$400$ random connected $n = 11,12$, and $622$ targeted hard-core-frame instances}
Lemma~\ref{lem:CAP} and Corollary~\ref{cor:CAP1} were tested on $1\,926$ $B$-universal low
vertices with $\degA = 1$ in every one and $0$ failures. Uniform random graphs almost never
land in the frame --- $12$ frame instances across the $2\,195$ exhaustive and random graphs
--- so a targeted generator was written: build $B$ with a chosen non-edge set, attach
$A$-vertices with random non-empty types, and keep only instances passing connected,
$\diam = 4$, non-forest, $f = \alpha+1$ and ``$A$ is a maximum independent set''. That
yields $622$ frame instances, of which $31$ satisfy $\Bloplus = \emptyset$; all five
conclusions of Theorem~\ref{thm:RIG} hold in $31$ of $31$. The scope demarcation of
Remark~\ref{rem:RIGscope} is measured, not assumed: on the hypothesis set the $(L,\nu)$
histogram is $\{(1,1):22,\ (1,3):2,\ (2,1):3,\ (3,1):4\}$, and the two instances with
$(L,\nu) = (1,3)$ have $\nu = 3 > L-1 = 0$ --- so the bound $\nu \le L-1$ is \emph{false} in
the frame and genuinely needs the reductio hypothesis, exactly as the proof says. None of
the $622$ frame instances satisfies $\res = \alpha$.
\end{remark}

\section{Statements at one adversarial review round}\label{sec:oneround}

Everything in this section is proved, and everything in it has passed exactly \emph{one}
review round --- a Qwen3.8-Max session ---
returning no mathematics defect against these statements. The second round, which the
campaign's protocol requires from a different model family before a statement is certified,
was dispatched and did not deliver: the brief cited four imports by name without stating
them, so the reviewing session verified every downstream calculation \emph{conditionally on
the quoted properties} and correctly refused to certify the imports it could not see. A
review round obstructed by a defect in its own brief is not a round. These statements are
therefore reported here in a separate tier, and the reader should treat them accordingly.

The setting is $\GFan(\tau,L,\nu)$ of Definition~\ref{def:fan} under the reductio
hypothesis; $C = \Blo \cup \{a_0\}$, $p = \tau - L = |\Bhi|$, $R = \sum_{a \in A'}\dg(a)$,
and $E$ is the escape count of the high phase as in Section~\ref{sec:fan}. The set-up
carries over verbatim from there: every $c \in C$ has $\dg(c) = \tau$ exactly, every
$w \in \Bhi$ has $\dg(w) \ge \tau+1$, every $a \in A'$ has $\dg(a) \le p < \tau$, and so
Lemma~\ref{lem:FAN1} applies word for word --- its proof uses only that the $\Bhi$ vertices
are the vertices of degree $\ge \tau+1$, and Lemma~\ref{lem:DICH}(c) for the rest.

\begin{lemma}[residue mass, general $\nu$]\label{lem:FAN4p}
\Sone{Qwen3.8-Max}
In $\GFan(\tau,L,\nu)$ under the reductio hypothesis, at the start of step $p+1$ the $A'$
entries sum to exactly $2\nu - E$, and the $C$-entries are $L + e_c$ for $c \in C$, where
$\sum_c e_c = E$.
\end{lemma}

\begin{proof}
Exactly the count of Lemma~\ref{lem:FAN4} with $\nu$ in place of $1$. By
Lemma~\ref{lem:FAN1} and Lemma~\ref{lem:DICH}(b),
$\sum_{j\le p} D_j = \sum_{\Bhi}\dg - \binom{p}{2}$. Here
$\sum_{x \in \Bhi}\degB(x) = p(\tau-1) - 2\nu$, since each of the $\nu$ non-edges lies
inside $\Bhi$ and is missed by both endpoints, and $\sum_{x\in\Bhi}\degA(x) = p + R$; so
$\sum_{\Bhi}\dg = p\tau - 2\nu + R$. Splitting the decrements by recipient ---
$\binom{p}{2}$ to later high heads, $p(L+1) - E$ to $C$, the rest to $A'$ --- and using
$p(L+1) = p\tau - p^2 + p$ and $2\binom{p}{2} = p^2 - p$ gives
$\mathrm{dec}_{A'} = R - 2\nu + E$, so the $A'$ total is $R - \mathrm{dec}_{A'} = 2\nu - E$.
Each $c \in C$ starts at $\tau$ and is decremented $p - e_c$ times, ending at $L + e_c$.
\end{proof}

\begin{lemma}[escape bound, general $\nu$]\label{lem:FAN8p}
\Sone{Qwen3.8-Max}
In $\GFan(\tau,L,\nu)$ under the reductio hypothesis, if $E \ge 1$ then $L \le 2\nu - E$.
\end{lemma}

\begin{proof}
Exactly the argument of Lemma~\ref{lem:FAN8}. At an escape step $t$ with $e_t \ge 1$
escapes, the head is high, so $D_t \ge \tau-t+2$; the $p-t$ remaining high vertices lie in
$\mathrm{block}_t$, and so do the $(L+1)-e_t$ non-escaping $C$-vertices; hence
$|\mathrm{block}_t \cap A'| \ge (\tau-t+2)-(p-t)-(L+1)+e_t = 1 + e_t \ge 2$, using
$\tau - p - L = 0$. Pick $x \in \mathrm{block}_t \cap A'$; as $\mathrm{block}_t$ is a prefix
and the escaping $c \notin \mathrm{block}_t$, $v_t(x) \ge v_t(c) \ge \tau-t+1$. By
Lemma~\ref{lem:FAN4p} the $A'$ mass at step $p+1$ is $2\nu-E$, so $v_{p+1}(x) \le 2\nu-E$;
and $x$ is not deleted in steps $t,\dots,p$, so $v_{p+1}(x) \ge v_t(x)-(p-t+1)$. Therefore
$\tau-t+1 \le (2\nu-E)+(p-t+1)$, i.e.\ $L = \tau-p \le 2\nu-E$.
\end{proof}

\begin{lemma}[backward induction, general form]\label{lem:FAN6p}
\Sone{Qwen3.8-Max}
In $\GFan(\tau,L,\nu)$ under the reductio hypothesis, the $A'$ multiset at the start of step
$p+1$ cannot have a unique maximum $w \ge 1$ whose second-largest entry is $\le w-2$.
\end{lemma}

\begin{proof}
As in Lemma~\ref{lem:FAN6}, with the block count
$a_j = D_j - (p-j) - (L+1) + e_j \ge 1 + e_j \ge 1$, so that a maximum entry of $M_j$ is
decremented at every step. Suppose $M_{j+1}$ has a unique maximum $v$ and no entry of value
$v-1$; the multiplicity argument of Lemma~\ref{lem:FAN6} shows $\mu = 1$ and that every
other entry of $M_j$ is at most $\max(\text{other entries of } M_{j+1})+1$. Hence, writing
$\sigma_j$ for the second-largest entry of $M_j$, the pair (unique maximum $v_j$,
$\sigma_j \le v_j - 2$) propagates backwards: from $(w,\sigma \le w-2)$ at $p+1$ one gets
$v_{p+1-t} = w+t$ and $\sigma_{p+1-t} \le \sigma + t \le v_{p+1-t}-2$ for every $t$, so the
induction never stops. At $t = p$ it gives $\max(M_1) = w+p$, while $M_1$ is the multiset of
$A'$ degrees and every $a \in A'$ has $N(a) \subseteq \Bhi$, so $\max(M_1) \le p$.
\end{proof}

Two conventions attach to Lemma~\ref{lem:FAN6p} and both were fixed by review findings.
First, for the purpose of the phrase ``second-largest entry'', a residue with a single
positive part is read as having second-largest $0$, whether or not a real zero entry is
present; the lemma's operative content is ``there is no entry equal to $w-1$'', and for a
single-part residue $[w]$ with $w \ge 2$ that is true under either reading, while $[1]$ is
excluded under either. Second, the residues the lemma kills are exactly those with a gap of
at least $2$ below a unique top: $[w]$ for every $w \ge 2$, $[3,1]$, $[4,1]$, $[4,2]$, and
so on; it does \emph{not} apply to $[1]$, $[1,1]$, $[2,1]$, $[2,2]$, $[2,1,1]$,
$[1,1,1,1]$. An earlier text wrote ``$[w]$ for every $w \ge 1$'', which is false at
$w = 1$ under the stated convention and undefined without it. No user is lost: every residue
this paper feeds to the lemma is $[2\nu]$ with $2\nu \ge 2$, or one of the survivors
tabulated in Appendix~\ref{app:enum}, all of which have $w \ge 3$.

\begin{theorem}[the $\nu = 2$ layer]\label{thm:GFAN2}
\Sone{Qwen3.8-Max}
For every $\tau$ and every $L \ge 3$ there is no graph carrying $\GFan(\tau,L,2)$ and
satisfying $\res(G) = \alpha(G)$.
\end{theorem}

\begin{proof}
Assume $\res = \alpha$. By Lemma~\ref{lem:FAN1} the heads of steps $1,\dots,p$ are $\Bhi$,
so exactly $L$ steps remain after step $p$; the multiset at the start of step $p+1$ is
$\{L+e_c : c \in C\}$ together with an $A'$ residue of total $2\nu - E = 4 - E$
(Lemma~\ref{lem:FAN4p}), and it must clear in exactly $L$ further steps.

\emph{Step 1: the escape budget.} By Lemma~\ref{lem:FAN8p}, $E \ge 1$ forces $L \le 4-E$.
So $L \ge 4$ gives $E = 0$, and $L = 3$ gives $E \le 1$.

\emph{Step 2: $L \ge 4$.} Here $E = 0$, the $C$-part is $[L]^{L+1}$ and the $A'$ residue is
a partition of $4$. The partitions $[4]$ and $[3,1]$ have a unique maximum with the next
entry at least $2$ below it, so Lemma~\ref{lem:FAN6p} kills both. For the remaining three,
run the process: while the common $C$-value is $t \ge 3$ the head is a $C$-entry, the other
$C$-entries are the only entries of value $\ge 3$, so the block is exactly those and the
$A'$ entries (all $\le 2$) are untouched; hence after $L-2$ steps the list is $[2]^3$
together with the untouched residue. Then
$[2,2]$: $[2]^5 \to [2,2,1,1] \to [1,1,0] \to 0$, three further steps, total $L+1$;
$[2,1,1]$: $[2,2,2,2,1,1] \to [2,1,1,1,1] \to [1,1,0,0] \to 0$, three further, total $L+1$;
$[1,1,1,1]$: $[2,2,2,1,1,1,1] \to [1]^6 \to \cdots$, four further, total $L+2$.
All exceed $L$, so none can occur.

\emph{Step 3: $L = 3$.} The complete list of admissible shapes is eight rows
($E \in \{0,1\}$; $E \ge 2$ is excluded by Step 1):
\begin{center}
\begin{tabular}{@{}cccl@{}}
\toprule
$E$ & $C$-part & $A'$ residue & steps to clear \\
\midrule
$0$ & $[3,3,3,3]$ & $[4]$ & $\mathbf{3}$, killed by Lemma~\ref{lem:FAN6p} \\
$0$ & $[3,3,3,3]$ & $[3,1]$ & $4 \ne 3$, and also killed \\
$0$ & $[3,3,3,3]$ & $[2,2]$ & $4 \ne 3$ \\
$0$ & $[3,3,3,3]$ & $[2,1,1]$ & $4 \ne 3$ \\
$0$ & $[3,3,3,3]$ & $[1,1,1,1]$ & $5 \ne 3$ \\
$1$ & $[4,3,3,3]$ & $[3]$ & $\mathbf{3}$, killed by Lemma~\ref{lem:FAN6p} \\
$1$ & $[4,3,3,3]$ & $[2,1]$ & $4 \ne 3$ \\
$1$ & $[4,3,3,3]$ & $[1,1,1]$ & $4 \ne 3$ \\
\bottomrule
\end{tabular}
\end{center}
Only the two single-part residues clear in exactly $L = 3$ steps, and
Lemma~\ref{lem:FAN6p} kills both. Hence no case survives.
\end{proof}

\begin{corollary}\label{cor:RIG1}
\Sone{Qwen3.8-Max}
In the hard core with $L = 2$, Proposition~\ref{prop:L2}(a),(b) force
$\Bloplus = \emptyset$; Theorem~\ref{thm:RIG} then gives Proposition~\ref{prop:L2}(c),(d)
directly, and $\nu \le L-1 = 1$ with $\nu \ge 1$ gives $\nu = \mbar = 1$, i.e.\
$\Fan(\tau,2)$.
\end{corollary}

This removes the tightness re-run from Proposition~\ref{prop:L2}(c): $\degA(b) = 1$ follows
from the degree cap alone. The campaign had pre-registered ``tightness arguments need the
inequality chain to be tight at every link, which is not spelled out'' as a likely failure
mode of Proposition~\ref{prop:L2} before the review round that examined it; the corollary is
what makes the concern moot.

\begin{corollary}\label{cor:GFAN2HC}
\Sone{Qwen3.8-Max}
In the hard core, if every low vertex is $B$-universal then $L \ge 4$ and $\nu \ge 3$.
\end{corollary}

\begin{proof}
Theorem~\ref{thm:RIG} makes the instance $\GFan(\tau,L,\nu)$ with $1 \le \nu \le L-1$;
$\nu = 1$ is $\Fan(\tau,L)$, killed for $L \ge 2$ by Theorem~\ref{thm:FAN}, and $\nu = 2$ is
killed for $L \ge 3$ by Theorem~\ref{thm:GFAN2}. $L = 1$ cannot occur, since
$1 \le \nu \le L-1$ already forces $L \ge 2$; and $L = 2$ forces $\nu = 1$. So $\nu \ge 3$,
whence $L \ge \nu+1 \ge 4$.
\end{proof}

\begin{corollary}\label{cor:GFAN2L3}
\Sone{Qwen3.8-Max}
In the hard core with $L = 3$, some low vertex is not $B$-universal
($\Bloplus \ne \emptyset$).
\end{corollary}

\begin{remark}[the other layer is not excluded at $L \ge 3$]\label{rem:otherlayer}
At $L = 2$, Proposition~\ref{prop:L2}(a),(b) forced $\Bloplus = \emptyset$ via
Theorem~\ref{thm:MB} and Fact~\ref{fact:Fb}. That argument does not survive $L = 3$, and we
record why so that no successor re-derives it as a theorem. With $k := |\Bloplus| = 1$ and
$\nu(\Blo) = 0$, Theorem~\ref{thm:MB} reads $1 + c + \mbar \le L - 1 = 2$, which is
satisfiable in two ways: $c = 0$, $\mbar = 1$, where $T_1,T_2 \subseteq \Bhi$ are singletons
spanning the unique $\Bhi$ non-edge and no contradiction with Fact~\ref{fact:Fb} arises
because $\mbar \ne 0$; and $c = 1$, $\mbar = 0$, where the non-universal low vertex is
itself in $T_1$ and $T_2$ sits inside its $\Bhi$ non-neighbourhood, so
Lemma~\ref{lem:Cstar} gives $\degA \ge 3$ there and Lemma~\ref{lem:CAP} forces $n_b \ge 2$
for that vertex. Both are budget-consistent. The contradiction that closed $L = 2$ used
$\mbar = 0$ \emph{and} $c = 0$ simultaneously, and at $L \ge 3$ the budget no longer forces
both. So the residual hard core splits as
\[
L \ge 3 \;=\; \bigl[\text{$B$-universal layer: } \GFan(\tau,L,\nu),\ 2 \le \nu \le L-1
\bigr] \;\cup\; \bigl[\Bloplus \ne \emptyset\bigr].
\]
Lemma~\ref{lem:CAP} is the handle on the second layer: $\degA(b) \le n_b+1$ converts every
$A$-degree lower bound into a non-edge lower bound. Concretely, Lemma~\ref{lem:Cstar} now
reads: every low vertex in $T_1 \cup T_2$ has at least two non-neighbours in $B$; and
Lemma~\ref{lem:pair}'s contribution to Theorem~\ref{thm:MB} is recovered by
Lemma~\ref{lem:CAP} for free.
\end{remark}

\section{A computer-assisted elimination, and a false quantifier}\label{sec:gfannu}

This section reports one theorem at the weakest tier used in this paper, and it reports it
together with the reason it sits there: the previous statement at this position was
\emph{false}, a review round produced a counterexample family, and the repaired statement
has not yet been reviewed. We consider the episode worth reporting in full rather than
silently repairing, because a pipeline that publishes only its successes gives no evidence
about its failure modes, and because the defect is of a species --- a quantifier admitting a
degenerate instance nobody constructed --- that no amount of numerical checking of the
non-degenerate range would have found.

\subsection{The statement, and the case list}

\begin{theorem}\label{thm:GFANnu}
\Szero{the statement below is the repaired one; the text that was reviewed was false, and
the repaired text is unreviewed}
For every $\tau$, every $\nu$ with $1 \le \nu \le 6$, and every $L \ge \nu+1$, there is no
graph carrying the configuration $\GFan(\tau,L,\nu)$ and satisfying $\res(G) = \alpha(G)$.
\end{theorem}

\begin{proof}
Four proved bounds make the case list finite for each $\nu$: Lemma~\ref{lem:FAN4p} (the
$A'$ residue totals $2\nu-E$ and the $C$-part is $\{L+e_c\}$ with $\sum e_c = E$);
Lemma~\ref{lem:FAN8p} ($E \ge 1 \Rightarrow L \le 2\nu-E$, so $E \ge 1$ occurs only for the
finitely many pairs $\nu+1 \le L \le 2\nu-E$, and $E \le \nu-1$); Corollary~\ref{cor:MB1}
($L \ge \nu+1$); and Lemma~\ref{lem:TAIL} (for $E = 0$ and $L \ge \lambda_1$, survival
depends on $\lambda$ alone, and the finitely many boundary rows
$\nu+1 \le L < \lambda_1 \le 2\nu$ are checked directly). Every shape in the resulting
finite list either fails to clear in exactly $L$ steps or is killed by
Lemma~\ref{lem:FAN6p}. The enumeration is reproduced in full in
Appendix~\ref{app:enum}; the counts are
\begin{center}
\begin{tabular}{@{}cccc@{}}
\toprule
$\nu$ & $E = 0$ rows Lemma~\ref{lem:FAN6p} misses & $E \ge 1$ shapes tested &
$E \ge 1$ rows it misses \\
\midrule
$1$ & none & $0$ & none \\
$2$ & none & $3$ & none \\
$3$ & none & $24$ & none \\
$4$ & none & $110$ & none \\
$5$ & none & $397$ & none \\
$6$ & none & $1\,211$ & none \\
\bottomrule
\end{tabular}
\end{center}
At $E = 0$ \emph{and} $L \ge \lambda_1$ --- inside Lemma~\ref{lem:TAIL}'s range, where
``clears in exactly $L$ steps'' is the $L$-free criterion $s_0(\lambda) = \lambda_1$ --- the
only residue clearing in exactly $L$ steps is the single part $[2\nu]$, for every
$1 \le \nu \le 6$, and $[2\nu]$ is killed by Lemma~\ref{lem:FAN6p}. Outside that range the
finitely many boundary rows $\nu+1 \le L < \lambda_1$ have their own survivors, tabulated in
Appendix~\ref{app:enum}, each with its own certificate.
\end{proof}

\begin{remark}[the demarcation, stated exactly]\label{rem:demarcation}
This is a computer-assisted proof and it is stated as such. The case list is finite
\emph{because} of Lemmas~\ref{lem:FAN4p}, \ref{lem:FAN8p}, \ref{lem:TAIL} and
Corollary~\ref{cor:MB1}, all hand-proved. Within that finite list, Appendix~\ref{app:enum}
prints the $E = 0$ criterion for all $159$ partitions of $2\nu$, $\nu \le 6$; all $51$
$E = 0$ boundary rows; the closed form
$S(\nu) = \sum_{E=1}^{\nu-1}(\nu-E)\,p(E)\,p(2\nu-E)$ reproducing the ``shapes tested''
column; and the complete roster of the $72$ surviving $E \ge 1$ rows with their
certificates. Every step except the \emph{completeness of the $E \ge 1$ survivor roster} is
therefore checkable by hand from the text alone; that one step is a seconds-scale machine
check over an explicitly specified, closed-form-counted finite set. $\nu = 1$ is
Theorem~\ref{thm:FAN} and $\nu = 2$ is Theorem~\ref{thm:GFAN2}, both with complete hand
proofs; $\nu = 3,\dots,6$ rest on the enumeration.
\end{remark}

An earlier version of the summary sentence in the proof read ``at $E = 0$ the only residue
clearing in exactly $L$ steps is the single part $[2\nu]$'', with no range qualifier, and
read literally that is contradicted by the paper's own data: $(L,\lambda) = (5,[7,1])$ at
$\nu = 4$, $(7,[9,1])$ at $\nu = 5$ and $(7,[10,1,1])$ at $\nu = 6$ are all $E = 0$, all
clear in exactly $L$ steps, and none is $[2\nu]$. They are boundary rows outside
Lemma~\ref{lem:TAIL}'s range, handled separately, and all are killed by
Lemma~\ref{lem:FAN6p}; the conclusion is unaffected, but the sentence as written was false,
and it was the punchline sentence. A review round found it, and the qualifier above is the
repair.

\begin{corollary}\label{cor:GFANnuHC}
\Szero{inherits Theorem~\ref{thm:GFANnu}'s state}
In the hard core, if every low vertex is $B$-universal then $\nu \ge 7$ and $L \ge 8$.
\end{corollary}

\begin{proof}
Theorem~\ref{thm:RIG} gives $\GFan(\tau,L,\nu)$ with $1 \le \nu \le L-1$ --- the lower bound
is clause (d), i.e.\ Observation~\ref{obs:R1} --- so $L \ge \nu+1$;
Theorem~\ref{thm:GFANnu} eliminates every $\nu$ with $1 \le \nu \le 6$, which by that lower
bound is every $\nu \le 6$ available here.
\end{proof}

\subsection{The refutation}\label{sec:refutation}

Theorem~\ref{thm:GFANnu}'s predecessor read ``for every $\nu \le 6$''. A review round
carried out by an OpenAI GPT-5.6 session through the \texttt{codex} interface returned the
verdict \textsc{refuted}, with the following observation: $\nu$ counts non-edges of $G[B]$,
so $\nu = 0$ is an admissible value, and the statement is false there.

\begin{proposition}[the counterexample family]\label{prop:K5}
\Scomp{$K_3, K_4, \dots, K_8$ all verified, every clause of Definition~\ref{def:fan}
checked individually, by two independent implementations}
Every complete graph $K_n$ with $n \ge 3$ carries the configuration
$\GFan(n-1,\,n-1,\,0)$ and satisfies $\res = \alpha$. Explicitly for $K_5$: it is connected
and not a forest, $\alpha = 1$ with $A = \{a_0\}$ a single vertex, $f = 2 = \alpha+1$, the
degree sequence is $[4,4,4,4,4]$ and $\res = 1 = \alpha$; $\tau = 4$, $\Blo = B$, $L = 4$,
$\Bhi = \emptyset$, $\nu = 0$, and $L \ge \nu+1$.
\end{proposition}

The proof of Theorem~\ref{thm:GFANnu} could never have covered $\nu = 0$, and inspecting why
is instructive. At $\nu = 0$, Lemma~\ref{lem:FAN4p} forces $E = 0$ and makes $\lambda$ the
\emph{empty} partition of $0$; Lemma~\ref{lem:TAIL} is stated for $\lambda$ non-empty; the
boundary branch needs $\lambda_1 \le 2\nu = 0$; and the $E \ge 1$ branch needs
$E \le \nu-1 = -1$. All three branches of the case split are vacuous at $\nu = 0$, which is
exactly why the enumeration in Appendix~\ref{app:enum} starts at $\nu = 1$ --- the printed
data already showed the gap, and nobody read it that way. The defect is a genuine gap
between statement and proof, not a rider; the repair is the missing quantifier
$1 \le \nu$, and nothing downstream changes, because every consumer reaches the theorem
through Theorem~\ref{thm:RIG}(d), where Observation~\ref{obs:R1} supplies $\nu \ge 1$.
Corollary~\ref{cor:GFANnuHC} is untouched for the same reason.

Three further things about the episode are worth stating, because they are the data a reader
should use to calibrate how much this pipeline's markers are worth.

\emph{First, the round was suspected of being shallow before it was read, and the suspicion
was wrong.} It ran faster and returned a smaller report than the preceding round, on a brief
carrying more named joints; the campaign recorded that as a prior for shallowness. Measured
against brief size the ranking inverts: the brief was $43$\,KB against the preceding round's
$84$\,KB and an earlier round's $138$\,KB, runtime tracked brief size almost linearly, and
per byte of brief this round produced more report than its predecessor. The lesson recorded
in the campaign's ledger is that predicting a reviewer's depth from a surface statistic
instead of from its work had by then failed three times.

\emph{Second, the round did substantially more than referee.} It built its own graph rather
than assuming the conclusion of Theorem~\ref{thm:RIG}: an eight-vertex frame instance with
$\alpha = 4$, $\diam = 4$, $f = 5 = \alpha+1$, connected, not a forest,
$\Blo = \{b_0,b_3\}$ both $B$-universal, $\Bhi = \{b_1,b_2\}$, unique $B$ non-edge
$b_1b_2$ so $\nu = 1$, all five conclusions of Theorem~\ref{thm:RIG} true, and
$\res = 3 \ne 4 = \alpha$. Every claim about it was reproduced independently. That last line
is what makes it a control rather than a witness: the frame holds and the reductio
hypothesis fails, so Lemmas~\ref{lem:FAN4p}, \ref{lem:FAN8p} and~\ref{lem:FAN6p} are
unavailable on it, and the report says so and does not use it as evidence for them.
It also declined to launder vacuity, stating in its own words that the positive-$\nu$
computation enumerates \emph{shapes}, not realized graphs, so that the absence of an
unkilled shape is not empirical evidence that no such graph exists.

\emph{Third, the campaign's own gate arithmetic went against the campaign.} Reading the
report's own joint-level verdicts, one could argue that the repaired statement
$1 \le \nu \le 6$ was reviewed in substance, since the specification, data and kill joints
all came back clean and one of them says ``clean for $1 \le \nu \le 6$''. That reading was
put to the orchestrator and \emph{rejected}, on the pre-registered rule that a clean round
counts only for the text as reviewed, and that a round which refutes a statement is not that
statement's clean round. Theorem~\ref{thm:GFANnu} is therefore recorded at zero rounds, not
at one, and Section~\ref{sec:open} names the single round that would close it.

\begin{remark}[the data were reproduced a third time]\label{rem:thirdtime}
\Scomp{Appendix~\ref{app:enum} recomputed in full from the specification alone by an
independent implementation; $0$ mismatches on the items listed below}
Because a refuting round raises the question of what \emph{else} is wrong, the whole of
Appendix~\ref{app:enum} was recomputed by an implementation written from scratch --- its own
Havel--Hakimi and residue primitives, its own independence-number and forest-number routines
by exhaustive subset enumeration, its own partition generator, its own clause checker for
Definition~\ref{def:fan} --- which reuses nothing from the reviewing session or its checker,
and which diffs its recomputation against the campaign's recorded data item by item, so that
a transcription slip would surface rather than be laundered. Agreement was exact on: the
$159$ values of $s_0(\lambda)$ and the fact that the printed partition set is exactly the
set of all $p(2\nu)$ partitions at each $\nu$; the $E = 0$ tail survivors
$[2],[4],[6],[8],[10],[12]$; the boundary pair counts $0,1,3,7,14,26$; the boundary
survivors and their certificates; $S(\nu) = 0,3,24,110,397,1211$ both by enumeration and by
closed form; the $E \ge 1$ survivor counts $0,1,4,9,20,38$, every row and every certificate,
with no row printed here that the recomputation lacks and none the other way; $0$ misses of
Lemma~\ref{lem:FAN6p} over all $9+72+6$ surviving rows; the $1\,817$ tail-formula pairs; and
the zero-padding inertness control, extended to $18\,963$ (list, padding) pairs with $0$
padding-dependent step counts.
\end{remark}

\section{What remains}\label{sec:open}

\subsection{The residual configuration space}

Collecting Theorems~\ref{thm:tau2}, \ref{thm:T3}, \ref{thm:FAN} and~\ref{thm:RIG} with
Corollaries~\ref{cor:FANHC} and~\ref{cor:RIG2}, a counterexample to the diameter-four
dichotomy \eqref{eq:D4} --- equivalently, a graph in the hard core of
Definition~\ref{def:frame} --- must satisfy
\[
\tau \ge 4, \qquad L \ge 3,
\]
and must lie in one of the two layers of Remark~\ref{rem:otherlayer}: either every low
vertex is $B$-universal, in which case Theorem~\ref{thm:RIG} pins the graph to
$\GFan(\tau,L,\nu)$ with $2 \le \nu \le L-1$ (and, granting Theorem~\ref{thm:GFANnu} at its
current tier, with $\nu \ge 7$ and $L \ge 8$); or some low vertex is not $B$-universal, in
which case no rigidity statement is available and Lemma~\ref{lem:CAP} is the only handle
so far --- it converts every $A$-degree lower bound into a lower bound on the number of
non-neighbours in $B$, which is what a hypothetical $L$-uniform argument would consume.

Nothing here bounds $\tau$ or $L$ from above, and there is a measured reason to expect no
such bound: the family with $B$ independent, $c$ universal $A$-vertices and two private
leaves per $B$-vertex lands genuinely inside the hard-core frame for every $\tau \ge 3$ and
$c \ge \tau$, so a case analysis proceeding one $\tau$ at a time can never terminate. That
observation is why every statement in Sections~\ref{sec:toolkit}--\ref{sec:rig} is
$\tau$-uniform.

\subsection{Two combinatorial conjectures that would close the $B$-universal layer}

The table underlying Theorem~\ref{thm:GFANnu} is uniform in $\nu$ as far as it was computed:
at $E = 0$ only $\lambda = [2\nu]$ survives, and every $E \ge 1$ survivor has a unique
maximum with a gap of at least $2$ below it. If that persists for all $\nu$, the entire
$B$-universal layer of the hard core is empty at every $L$, and by
Remark~\ref{rem:otherlayer} the whole hard core reduces to the layer
$\Bloplus \ne \emptyset$. Precisely:

\begin{conjecture}\label{conj:C1}
\Sconj{verified for every $\nu \le 10$; proved for $\lambda = [2\nu]$ and for every
$\lambda$ with at most two parts}
For every $\nu$ and every partition $\lambda$ of $2\nu$ other than $[2\nu]$,
$s_0(\lambda) \ne \lambda_1$.
\end{conjecture}

\begin{conjecture}\label{conj:C2}
\Sconj{verified for every $\nu \le 10$}
Every $(L,E)$ survivor with $E \ge 1$ has an $A'$ residue $\lambda$ with a unique maximum
whose second-largest entry is at most $\lambda_1 - 2$.
\end{conjecture}

\begin{remark}[partial progress on Conjecture~\ref{conj:C1}]\label{rem:C1progress}
\Szero{hand proofs; never queued for review}
Four statements were proved by hand and are recorded here at the weakest tier, having never
been submitted to any review round. (i) $s_0([2\nu]) = 2\nu$ for every $\nu$, by a two-step
parity induction with no computation --- so the $[2\nu]$ row of the enumeration is
hand-proved at every $\nu$, not only for $\nu \le 6$. (ii) Conjecture~\ref{conj:C1} is not a
step-count statement but a residue statement: writing $k$ for the number of parts of
$\lambda$, one has $s_0(\lambda) = \lambda_1$ if and only if
$\res\bigl([\lambda_1]^{\lambda_1+1} \cup \lambda\bigr) = k+1$. (iii) Consequently, by
Fact~\ref{fact:fms}, Conjecture~\ref{conj:C1} holds for a given $\lambda$ as soon as
$[\lambda_1]^{\lambda_1+1} \cup \lambda$ admits a realization with independence number at
most $k$ --- so the general case acquires a sufficient condition. That condition is
not known to be satisfiable in general: for $\lambda = (w,1)$ with $w \ge 3$ odd, no
realization with independence number at most $k$ exists, although $C1$ holds there.
It is therefore a sufficient condition and not a reduction. (iv) For $k \le 2$
the construction succeeds: $s_0\bigl((w,c)\bigr) = w+1$ exactly, whenever $w+c$ is even,
which it always is here since $\lambda \vdash 2\nu$; hence Conjecture~\ref{conj:C1} holds
for every partition with at most two parts. The step counts were verified over all $10\,268$
terminating cases with $n \le 28$, and the enumeration of Appendix~\ref{app:enum} was
separately extended to $\nu \le 10$ with the same outcome. That extension is deliberately
\emph{not} folded into Theorem~\ref{thm:GFANnu}'s statement: it was produced while review
rounds on the $\nu \le 6$ text were in flight, and desynchronising a theorem from the text a
reviewer is holding invalidates the round. When it is folded in it must be folded as
$1 \le \nu \le 10$, or the hole of Section~\ref{sec:refutation} re-opens.
\end{remark}

\subsection{The one review round that would close the certification surface}
\label{sec:resume}

The mathematics of this line is, in the campaign's judgement, complete as far as it goes;
what is not complete is its certification. Exactly seven statements sit at one round
(Section~\ref{sec:oneround}) and two at zero (Section~\ref{sec:gfannu}), and a single
correctly-briefed round by a non-Qwen family on the current repaired text would give the
seven their second family and the two their first. The obstruction to the round that was attempted
is documented above and is entirely on the dispatching side: the brief must \emph{paste}, in
full statement form, the four imports whose absence obstructed it --- Lemma~\ref{lem:FAN1},
which was not even in the assumable list; Lemma~\ref{lem:DICH} clause (b) explicitly, rather
than the lemma's name; Proposition~\ref{prop:L2} clauses (a) and (b) verbatim; and
Corollary~\ref{cor:L1short}, or better, drop its citation entirely, since
Corollary~\ref{cor:GFAN2HC} no longer uses it. That round was specified but not dispatched,
the campaign having entered a wind-down.

\section{Limitations}\label{sec:limits}

We state the limitations of this paper plainly.

\begin{enumerate}
\item \emph{The conjecture is not resolved, and the gap is not small.} What is proved is a
reduction and a narrowing: the diameter-four dichotomy \eqref{eq:D4} implies the conjecture
at diameter four, and the configurations that could violate \eqref{eq:D4} have been narrowed
to $\tau \ge 4$, $L \ge 3$ and two layers, one of which is pinned to an explicit family and
one of which is not pinned at all. Nothing here bounds $\tau$ or $L$, and
Section~\ref{sec:open} explains why no argument proceeding one $\tau$ at a time can.

\item \emph{Nothing is formalized.} The companion papers of this pipeline
\cite{batch1,fernandespaper} end in Lean~4 files whose compilation is the arbiter of every
claim. This one does not, and the strongest marker used here is a statement about a review
process. Two independent reviewing sessions from different model families returning no
mathematics defect is materially weaker evidence than a kernel-checked proof, and the reader
should treat it as such. The one thing the tiering does guarantee is that the paper
distinguishes the levels rather than averaging them.

\item \emph{The residual is at diameter four only in a measured range.}
Corollary~\ref{cor:B2} settles $d \in \{1,2,3,5,6,9\}$ unconditionally; $d = 4$, $d = 7$,
$d = 8$ and every $d \ge 10$ are not settled by the elementary package. What
Remark~\ref{rem:coverage} establishes is that in the exhaustively checked range $n \le 8$
every graph missed by the three certificates has $d = 4$. Proposition~\ref{prop:dichotomy}
is therefore a statement about the diameter-four case and not about ``the residual case''
in general, and the campaign's own internal summary of it as the latter is an overstatement
that we correct here.

\item \emph{One certified proof is summarised, not reproduced.}
Theorem~\ref{thm:T3}'s complete case analysis --- four branches, several sub-branches, and
seven landed repairs from three review rounds --- runs to several pages of dense case work
in the source report, and Section~\ref{sec:smalltau} gives its skeleton and its tools rather
than every branch. The two review rounds certified the full text, not this summary. A reader
who wants the branch-level detail must go to the source report; a reader who wants to check
Theorem~\ref{thm:T3} independently is better served by the parameterisation in
Remark~\ref{rem:T3num}, which reduces the $\tau = 3$ frame to an exhaustive scan over a
box of multiplicities.

\item \emph{Fact~\ref{fact:fms} is used as a black box and its bibliographic data comes from
an automated pass.} The inequality $\res \le \alpha$ carries every argument in this paper.
It is not reproved here, and the attribution should be re-checked against the literature
before any formalization. It is asserted by the verification script on the whole exhaustive
range and never fails there, which is evidence about the statement and none at all about the
attribution.

\item \emph{Numerical scopes are exactly what the markers say.} Every
\Scomp{\dots} marker names a finite range, and a failed bounded search inside it is not a
proof of impossibility. In particular the hard core is empty on everything reachable at the
sizes searched --- which is what the conjecture predicts, and which therefore supplies no
discriminating evidence about the unproved statements at all. The mechanism checks reported
in Observation~\ref{obs:FANE} and elsewhere run \emph{off} the reductio hypothesis for
exactly this reason, and they corroborate the shape of an argument rather than its
conclusion.

\item \emph{Two of the campaign's own statements were false.} Remark~\ref{sec:upstream-selfref}
and Section~\ref{sec:refutation} report them, with their counterexamples. Both were caught
by the review gate, not by the solver. The honest reading of that is not that the gate
works, but that the solver's first-pass output should not be trusted without it; the
success rate of the first gate, not of the pipeline, is the measure of a solver.

\item \emph{Literature search was automated.} We checked for prior work on
Conjecture~\ref{conj:main} by automated search and found none beyond the collection and the
formal corpus, but automated searches miss things. No priority is claimed for anything here,
and readers should assume that some of the elementary material in
Section~\ref{sec:elementary} is known.

\item \emph{Seven statements are one round short and two are two rounds short.} They are
segregated into Sections~\ref{sec:oneround} and~\ref{sec:gfannu} and marked accordingly.
Corollary~\ref{cor:GFANnuHC} in particular should not be quoted as established; the
strongest unconditional consequence of the certified material is
Corollary~\ref{cor:FANHC}'s $L \ge 3$.
\end{enumerate}

\subsection*{Acknowledgement of process}

No human contributed a mathematical step to this paper. The pipeline --- problem selection,
attack design, solving, adjudication, adversarial review, computation and drafting --- was
run autonomously. Responsibility for the correctness of what is claimed nevertheless rests
where it always does: with the reader who checks it, and, for the statements below the top
tier, with the reader who does not yet.

\appendix

\section{Certification ledger}\label{app:ledger}

Table~\ref{tab:ledger} names, for each statement of this paper carrying a
\textsf{certified} or \textsf{one round} marker, the campaign's registry row and the review
rounds behind it. ``Family'' records the model family of each round; the protocol
requires the two rounds of a certified statement to come from different families, and
requires at least one of them to have read the final repaired text. The adjudication column
names the campaign's own reproduction script, which re-derives every claim of the report
from scratch before it is upheld; those scripts are listed in
Appendix~\ref{app:scripts}.

\begin{center}
\small
\begin{longtable}{@{}p{0.055\textwidth}p{0.30\textwidth}p{0.30\textwidth}p{0.25\textwidth}@{}}
\toprule
row & statements & review rounds (family) & adjudication \\
\midrule
\endfirsthead
\toprule
row & statements & review rounds (family) & adjudication \\
\midrule
\endhead
R-1 & Lemma~\ref{lem:Zplus} & round A2 (Qwen3.8-Max) + round B (Claude Opus 5) +
confirmation pass (Claude Opus 5, separate session) &
\texttt{w61\_r4\_a2check.py}, \texttt{w61\_thmK\_check.py} \\
R-2 & Lemma~\ref{lem:DICH} & same chain & same \\
R-3 & Theorem~\ref{thm:K} & same chain & same \\
R-4 & Corollary~\ref{cor:K1} & same chain & same \\
R-5 & Lemma~\ref{lem:F3prime} & same chain & \texttt{w61\_adjudicate\_r4.py} \\
R-6, R-7 & the two superseded counting statements of Remark~\ref{sec:upstream-selfref} &
round B (Claude Opus 5) + confirmation pass & \texttt{w61\_adjudicate.py} \\
R-9 & Lemmas~\ref{lem:FAN1}, \ref{lem:FAN3}, \ref{lem:FAN4}, \ref{lem:FAN6},
\ref{lem:FAN7}, \ref{lem:FAN8} &
rounds A, A2, B (Claude Opus 5, three separate sessions) + cross-family round
(Qwen3.8-Max) + verification pass (Qwen3.8-Max) &
\texttt{w61\_r5\_a2check.py}, \texttt{w61\_r6\_q14check.py} \\
R-10 & Theorem~\ref{thm:FAN} & same chain & same \\
R-13 & Theorem~\ref{thm:T3} with repairs R1, R3, J1, J2, K1, K2, K3 &
round A (Claude Opus 5) + Part 2 of the verification pass (Qwen3.8-Max), both
math joints re-derived independently &
\texttt{w61\_r5\_jcheck.py}, \texttt{w61\_tau3\_struct.py} \\
R-14 & Observation~\ref{obs:FAN5} & released by R-15; own content certified with R-9 &
\texttt{w61\_r5\_bcheck.py} \\
R-15 & Lemma~\ref{lem:TAIL} with repairs Y1--Y3 &
round 1 (Qwen3.8-Max, no mathematics defect) + round 2 (\texttt{codex} \texttt{spark},
clean on the repaired text) & \texttt{w61\_r11\_adjudicate.py},
\texttt{w61\_S3\_TAIL\_spark\_check.py} \\
R-16 & Theorem~\ref{thm:LOW}, Theorem~\ref{thm:SL}, Lemma~\ref{lem:HI} &
round 1 (Qwen3.8-Max, no mathematics defect, repairs X1--X4) + round 2
(\texttt{codex} \texttt{sol}, on the repaired text, Steps 1--6 of
Theorem~\ref{thm:SL} re-derived independently) & \texttt{w61\_r10\_adjudicate.py},
\texttt{w61\_r12\_adjudicate.py} \\
R-17 & Theorem~\ref{thm:MB}, Corollaries~\ref{cor:L1short}, \ref{cor:MB1},
\ref{cor:L1prime}, Propositions~\ref{prop:L1}, \ref{prop:L2} & same two-family chain;
all joints clean in round 2 with independent re-derivation, including the two-sided
squeeze of Proposition~\ref{prop:L2}(c) & same \\
R-18 & Corollary~\ref{cor:SLHC} with repair AA1 & same chain; AA1's replacement argument
certified directly by the orchestrator & \texttt{w61\_r12\_adjudicate.py} \\
R-19 & Corollary~\ref{cor:FANHC} & released by R-16--R-18; own content certified with
R-9/R-10 & --- \\
R-20 & Lemma~\ref{lem:CAP}, Corollary~\ref{cor:CAP1} with repair AB3 &
round 1 (Qwen3.8-Max) + round 2 (\texttt{codex} \texttt{sol}; both bounds re-derived
independently) & \texttt{w61\_r13\_adjudicate.py} \\
R-21 & Theorem~\ref{thm:RIG}, Corollary~\ref{cor:RIG2} with repair AB2 &
same two families; round 2 supplied an independently constructed frame witness
(Section~\ref{sec:refutation}) & \texttt{w61\_r13\_adjudicate.py} \\
\midrule
--- & Lemmas~\ref{lem:FAN4p}--\ref{lem:FAN6p}, Theorem~\ref{thm:GFAN2},
Corollaries~\ref{cor:RIG1}, \ref{cor:GFAN2HC}, \ref{cor:GFAN2L3} &
\textbf{one round only}: Qwen3.8-Max, joints clean, repairs Z1--Z4 landed. The second
family's round was obstructed by a briefing defect (Section~\ref{sec:oneround}) &
\texttt{w61\_r12\_adjudicate.py} \\
--- & Theorem~\ref{thm:GFANnu}, Corollary~\ref{cor:GFANnuHC} &
\textbf{no completed round}: the text reviewed was false
(Section~\ref{sec:refutation}); the repaired text is unreviewed &
\texttt{w61\_r13\_adjudicate.py} \\
\bottomrule
\caption{Certification ledger. Registry rows are the campaign's own identifiers.}
\label{tab:ledger}
\end{longtable}
\end{center}

Two structural facts about this table are worth stating. First, the model families are
genuinely separated: no reviewing session read another's report, the campaign's dispatch
records, or the pre-declared failure modes, and each brief carried an explicit isolation
rule. Second, three rows were held back at points where accepting a clean verdict wholesale
would have promoted something unearned --- Observation~\ref{obs:FAN5} was held until
Lemma~\ref{lem:TAIL} cleared its own two-family bar, Corollary~\ref{cor:FANHC} was held
until the whole of Section~\ref{sec:budget} cleared, and the entire
Section~\ref{sec:oneround} block was held and remains held. Those holds are the part of the
ledger that carries information.

\section{Verification-artifact inventory}\label{app:scripts}

The scripts below are Python programs in the campaign repository; none requires a solver,
and the longest single run is about $100$ seconds. Most are stand-alone. Three classes are
not: \texttt{w61\_r6\_q14check.py}, \texttt{w61\_r13\_adjudicate.py} and the four brief
builders read the campaign's working notes at run time and so run only inside the
repository, and \texttt{w61\_S3B\_T3\_probe.py} is a library of probe routines with no
executable main.

Eight of the scripts --- \texttt{w61\_S3\_76\_sol\_check.py},
\texttt{w61\_S3\_GFAN\_sol\_check.py}, \texttt{w61\_S3\_TAIL\_spark\_check.py},
\texttt{w61\_r6\_q14check.py}, and \texttt{w61\_r10\_adjudicate.py} through
\texttt{w61\_r13\_adjudicate.py} --- begin by recomputing the two calibration values
$\res(K_2) = 1$ and $\res(C_n) = \lceil n/3 \rceil$ for $n = 3,\dots,9$ from their own
implementation, so that a mistranscription of the residue definition aborts the run before
any reported number is produced. The others do not. Paths are relative to the repository
root.

\subsection*{B.1 The falsification and certificate stages}

\begin{description}[leftmargin=1.5em,style=nextline]
\item[\texttt{notes/proofs/wowii61\_verify.py}]
The primary verifier. Mode \texttt{all} runs stages A--E of Remark~\ref{rem:falsify} and
\emph{asserts} $\res \le \alpha$ (Fact~\ref{fact:fms}), $f \ge \alpha+1$
(Lemma~\ref{lem:f-alpha}), $f \ge \alpha + \lfloor (d-1)/4 \rfloor + 1$
(Corollary~\ref{cor:B1}) and \eqref{eq:C61} itself on every graph it touches, so a
counterexample to any of them aborts the run. Mode \texttt{w} emits the witnesses
$W_1,\dots,W_7$ of Section~\ref{sec:obstacles} with all their invariants. Mode \texttt{f}
evaluates the two sub-conjectures of the campaign's first round and the two dead routes of
Section~\ref{sec:obstacles} over $46\,995$ graphs. It also produces the coverage table of
Remark~\ref{rem:coverage}. \texttt{networkx} is used only for the atlas of graphs on at most
seven vertices; everything else is bit manipulation.
\end{description}

\subsection*{B.2 Toolkit and small-$\tau$}

\begin{description}[leftmargin=1.5em,style=nextline]
\item[\texttt{problems/wowii/w61\_lemH.py}]
Certifies Lemmas~\ref{lem:S} and~\ref{lem:F3prime} and the head-decay bound they replaced,
over $116\,873$ connected graphs ($n \le 7$ exhaustive, the $n = 8$ cover, $12\,000$ random
$n = 9,\dots,12$, structured families), of which $36\,612$ satisfy the reductio hypothesis;
canonical sort plus three randomised adversarial tie-breaks per graph.
\item[\texttt{problems/wowii/w61\_blockocc.py}]
Certifies Lemma~\ref{lem:Zplus} and the block-occupancy mechanism over $116\,189$ connected
graphs, $36\,664$ of them under the reductio hypothesis, with the six labelled claims
reported separately.
\item[\texttt{problems/wowii/w61\_thmK\_check.py}]
Independent check of Theorem~\ref{thm:K} (Remark~\ref{rem:thmK}), written by the owner
agent and reusing no code from the agent that proposed the theorem: $114\,914$ connected
graphs, $1\,464$ hypothesis hits, $0$ non-cliques.
\item[\texttt{problems/wowii/w61\_cstar.py}]
Certifies Lemma~\ref{lem:Cstar} on the hard-core \emph{frame} --- no reductio hypothesis
imposed --- over the exhaustive $n \le 7$ atlas plus the $n = 8$ cover: $332$ frame graphs,
$399$ pairs realising Fact~\ref{fact:Fb}, $0$ failures, plus the $\tau$-histogram of
Remark~\ref{rem:obstacle-tau}.
\item[\texttt{problems/wowii/w61\_famgen.py}]
Generates the two unbounded hard-core-frame families used in Section~\ref{sec:open} and
confirms membership instance by instance for $\tau = 3,\dots,7$ and $c = 2,\dots,10$.
\item[\texttt{problems/wowii/w61\_tau3.py}]
Exhaustive scan of the $\tau = 3$ frame by the parameterisation of
Remark~\ref{rem:T3num}: $1\,119\,744$ parameter tuples in a multiplicity box of size $5$,
of which $591\,710$ lie in the frame; reports the $e_B$ distribution and the minimum of
$m - \sum_{i \le 3} D_i$.
\item[\texttt{problems/wowii/w61\_tau3\_struct.py}]
Certifies the structural skeleton of Theorem~\ref{thm:T3} --- the three $B$-degree bounds
and the two implications $e_B = 1 \Rightarrow k \ge 1$, $e_B = 2 \Rightarrow k \ge 2$ ---
on $147\,348$ instances, and runs the four explicit families of the proof at parameters up
to $X = 59$.
\item[\texttt{problems/wowii/w61\_r5\_t3check.py}]
Corroborating box for Theorem~\ref{thm:T3}: $12\,288$ graphs with $\tau = 3$ over all $e_B$
and all type multiplicities $\le 2$, with $A$ verified maximum; $0$ hard-core instances.
\item[\texttt{problems/wowii/w61\_r3\_repair.py}, \texttt{w61\_r5\_jcheck.py}]
Numerical backing for the repairs to Theorem~\ref{thm:T3}, including the direct
Havel--Hakimi computation on the unique $k = 0$ graph and the refutation of the enumeration
claim that repair K2 replaced by a closed-form diameter argument.
\end{description}

\subsection*{B.3 The budget chain}

\begin{description}[leftmargin=1.5em,style=nextline]
\item[\texttt{problems/wowii/w61\_r4\_low.py}]
Certifies \eqref{eq:LOW1}--\eqref{eq:LOW3}, the position-independence identity in the proof
of Theorem~\ref{thm:LOW}, both clauses of Lemma~\ref{lem:DICH}, Lemma~\ref{lem:HI}, and
Theorem~\ref{thm:K} as the case $L = 0$: $6\,403$ connected graphs, $34\,040$ reductio
trajectory checks, every maximum independent set tested rather than one, canonical sort plus
three adversarial tie-breaks. Emits the per-$L$ slack minima quoted in
Remark~\ref{rem:SLsharp} and the hard-core-frame $L$ histogram.
\item[\texttt{problems/wowii/w61\_r4\_slack.py}]
The corpus deliberately enriched with near-split graphs ($B$ a clique minus a few random
edges), which is the regime Proposition~\ref{prop:L1} lives in: $5\,913$ graphs,
$26\,514$ checks, $1\,240$ instances with $L = 1$, in all of which
$(\degA(b_0),\mbar) = (1,0)$.
\item[\texttt{problems/wowii/w61\_r4\_thmSL.py}]
Certifies \emph{every intermediate assertion} of the proof of Theorem~\ref{thm:SL}
separately, not only its conclusion: the eleven labelled steps, over $22\,450$ trajectories
with five tie-breaks each. Produces the $s_j$ histogram of Remark~\ref{rem:SLsharp} and the
$1\,925$ counterexamples to the stronger statement the theorem is \emph{not}.
\item[\texttt{problems/wowii/w61\_r4\_mb.py}]
Numerical backing for Theorem~\ref{thm:MB} and Propositions~\ref{prop:L1},
\ref{prop:L2}.
\item[\texttt{problems/wowii/w61\_r5\_generalS.py}]
Closes the general-$s$ species: checks each toolkit lemma off the reductio hypothesis, in
the regime where a reductio-only statement would be invalid.
\end{description}

\subsection*{B.4 The fan family and the rigidity chain}

\begin{description}[leftmargin=1.5em,style=nextline]
\item[\texttt{problems/wowii/w61\_r4\_fanL.py}]
Enumerates $\Fan(\tau,L\ge2)$ degree sequences over a parameter box and reports
$\alpha - \res$ for each. In its default box $\tau \le 9$, $R \le 20$, $|A'| \le 16$:
$736\,630$ strict and $1\,101\,234$ superset sequences, $0$ with $\res = \alpha$, and the
control $L = 1$ with $1\,196\,636$ sequences and $0$ survivors. The enlarged box
$\tau \le 10$, $R \le 20$, $|A'| \le 14$, invoked as
\texttt{w61\_r4\_fanL.py 10 20 14}, reaches $1\,098\,141$ strict and
$1\,700\,094$ superset sequences; this is the source of Observation~\ref{obs:FANE}.
\item[\texttt{problems/wowii/w61\_r4\_fanmech.py}]
The mechanism check of Section~\ref{sec:conventions}'s firewall: $64\,911$ labelled
Havel--Hakimi runs under canonical and randomised tie-breaks, testing which entries occupy
which blocks. Includes a counterfactual-availability control --- the same test applied to
shapes breaking exactly one hypothesis of Definition~\ref{def:fan} returns
$\res = \alpha$ in $36\,650$ of $46\,138$ $B$-clique cases, so the test discriminates and is
not vacuously negative.
\item[\texttt{problems/wowii/w61\_r4\_fanres.py}, \texttt{w61\_r4\_fanE.py},
\texttt{w61\_r4\_fan.py}]
Certify the residue-mass count of Lemma~\ref{lem:FAN4}, the five structural hypotheses of
Lemma~\ref{lem:FAN6} ($508\,239$ high-phase runs, three tie-breaks each) and the escape
budget of Lemma~\ref{lem:FAN7}; report the residue partition in $450\,303$ superset
sequences.
\item[\texttt{problems/wowii/w61\_r5\_bcheck.py}]
Independent confirmation of Observation~\ref{obs:FAN5}'s second half for
$L = 2,\dots,9$ against direct simulation.
\item[\texttt{problems/wowii/w61\_r5\_rig.py}]
Certifies Lemma~\ref{lem:CAP}, Corollary~\ref{cor:CAP1} and all five conclusions of
Theorem~\ref{thm:RIG}, with its own labelled residue implementation, its own graph code and
its own corpora, importing no earlier script. Contains the targeted frame generator of
Remark~\ref{rem:RIGnum} and the $(L,\nu)$ histogram that measures the scope demarcation of
Remark~\ref{rem:RIGscope}.
\item[\texttt{problems/wowii/w61\_r5\_gfan2.py}]
Certifies Theorem~\ref{thm:GFAN2}'s three trajectories and its eight-row $L = 3$ table, and
runs the $\nu \le 6$ enumeration through the three functions
\texttt{tail\_check} / \texttt{escape\_check} / \texttt{complete\_check}.
\item[\texttt{problems/wowii/w61\_r12\_gfannu\_embed.py}]
The instrumented re-run that produced Appendix~\ref{app:enum}: the same check as the
previous entry, re-instrumented so that every intermediate column is printed rather than
summarised, together with the zero-padding inertness control and the closed-form evaluation
of $S(\nu)$ independent of the enumeration loop.
\item[\texttt{problems/wowii/w61\_r13\_padding\_ext.py}]
Extends the padding-inertness control to full scope: $964$ rows with $E = 0$ and all
$1\,745$ enumerated $E \ge 1$ shapes for $\nu \le 6$, times seven paddings, giving
$18\,963$ (list, padding) pairs with $0$ padding-dependent step counts.
\item[\texttt{problems/wowii/w61\_r12\_gfannu\_ext.py}]
The $\nu \le 10$ extension of Remark~\ref{rem:C1progress}, to the same specification.
\item[\texttt{problems/wowii/w61\_r12\_c1\_probe.py}, \texttt{w61\_r12\_c1k2.py}]
Probe and certificate for Conjecture~\ref{conj:C1}: the first evaluates $s_0(\lambda)$ over
every partition of every relevant weight in range; the second certifies the two-part case
over all $10\,268$ terminating instances with $n \le 28$.
\end{description}

\subsection*{B.5 Adjudication scripts}

Each review round is reproduced from scratch before any of its findings is upheld. These
scripts are written to be maximally hostile to the campaign's own working text: they
recompute the reported data from the specification alone, with independent implementations,
and diff the result against the campaign's recorded data, so that a transcription slip
surfaces rather than being laundered.

\begin{description}[leftmargin=1.5em,style=nextline]
\item[\texttt{problems/wowii/w61\_adjudicate.py}, \texttt{w61\_adjudicate\_r4.py},
\texttt{w61\_r4\_a2check.py}, \texttt{w61\_r5\_a2check.py},
\texttt{w61\_r6\_q14check.py}]
Reproductions of the rounds behind rows R-1--R-13, each with an independent labelled
Havel--Hakimi implementation.
\item[\texttt{problems/wowii/w61\_r10\_adjudicate.py},
\texttt{w61\_r11\_adjudicate.py}, \texttt{w61\_r12\_adjudicate.py},
\texttt{w61\_r13\_adjudicate.py}]
Reproductions of the four later rounds: respectively the budget chain, Lemma~\ref{lem:TAIL},
the second-family pass on Sections~\ref{sec:budget}--\ref{sec:fan}, and the refuting round
of Section~\ref{sec:refutation}. The last is the one described in
Remark~\ref{rem:thirdtime}; it reports \texttt{FAILURE\_COUNT = 0}.
\item[\texttt{problems/wowii/w61\_S3\_76\_sol\_check.py},
\texttt{w61\_S3\_GFAN\_sol\_check.py}, \texttt{w61\_S3\_TAIL\_spark\_check.py},
\texttt{w61\_S3B\_K\_hh.py}, \texttt{w61\_S3B\_K\_exh.py},
\texttt{w61\_S3B\_K\_avail.py}, \texttt{w61\_S3B\_K\_probe.py},
\texttt{w61\_S3B\_T3\_probe.py}]
The reviewing sessions' \emph{own} checkers, archived as they were received. They were
re-run independently during adjudication; one of them --- the round-B judge's exhaustive
enumerator --- performed $5\,842\,009$ trajectories under exhaustive tie-break enumeration
over $28\,163$ graphs.
\item[\texttt{problems/wowii/w61\_r10\_build\_spark\_brief.py},
\texttt{w61\_r10\_build\_sol\_brief.py},
\texttt{w61\_r12\_build\_gfannu\_brief.py},
\texttt{w61\_r12\_build\_sol\_gfan\_brief.py}]
The brief builders. Each performs a leak check before emitting, verifying that the brief
contains no expected answer, no pre-declared failure mode and no other reviewer's report.
The defect described in Section~\ref{sec:oneround} is an omission in the last of these ---
four statements cited by name and not pasted --- and repairing it is a change to that
script's appendix block.
\end{description}

\section{The \texorpdfstring{$\nu \le 6$}{nu <= 6} enumeration, in full}\label{app:enum}

This appendix is what makes Theorem~\ref{thm:GFANnu} reproducible from this text rather
than from an archived file. Its predecessor cited the enumeration by filename, and a
reviewing session correctly refused to certify a column it could not recompute; everything
below is the repair.

\subsection*{C.1 The object being enumerated, specified exactly}

At the start of step $p+1$ of a $\GFan(\tau,L,\nu)$ run under the reductio hypothesis, the
Havel--Hakimi value list is, by Lemma~\ref{lem:FAN4p}, exactly: a $C$-part of $L+1$ entries,
the $c$-th equal to $L + e_c$ with $e_c \ge 0$ and $\sum_c e_c = E$; an $A'$-part which is a
partition $\lambda$ of $2\nu - E$; and zero entries. Five conventions fix the enumeration.

\begin{enumerate}
\item \emph{Unlabelled.} The step count of a Havel--Hakimi list depends only on the value
multiset, so the $C$-part is enumerated as a multiset $\{L+e_c\}$ --- i.e.\ $e$ is a
partition of $E$ into at most $L+1$ non-negative parts --- and the $A'$-part as a partition
of $2\nu-E$. Distinct labellings of the same multiset are the same row.

\item \emph{Zero entries are inert}, verified rather than assumed. The block at each step is
the $D$ largest non-head entries, so a zero enters the block only when fewer than $D$
positive entries remain, and in that case the run aborts either way; the zero padding in the
scripts is therefore a convenience and $|A'|$ never has to be pinned down. As a control, all
$964$ rows with $E = 0$ and all $1\,745$ enumerated $E \ge 1$ shapes for $\nu \le 6$ were
padded with $0,1,2,3,4,8,13$ extra zeros --- $18\,963$ (list, padding) pairs --- with $0$
padding-dependent step counts.

\item \emph{No graphicality filter is applied, and none is needed.} The enumeration is a
superset argument: it ranges over every multiset the reductio hypothesis \emph{could}
produce and shows that none of them clears in exactly $L$ steps except rows that
Lemma~\ref{lem:FAN6p} independently forbids. Adding a graphicality or feasibility filter can
only remove rows, so the conclusion is monotone in the right direction.

\item \emph{``Clears in exactly $L$ steps''} means: iterating
\emph{head-deletes-the-$D$-largest} from the list above reaches all-zeros in exactly $L$
deletions. A run that would drive a zero entry negative, or whose head exceeds the number of
remaining entries, returns ``not a step sequence'' and is not a survivor.

\item \emph{Parameter ranges.} $L \ge \nu+1$ by Corollary~\ref{cor:MB1};
$E \ge 1 \Rightarrow L \le 2\nu-E$ by Lemma~\ref{lem:FAN8p}, which together force
$E \le \nu-1$; and the $E = 0$ rows split at $\lambda_1$ into the range $L \ge \lambda_1$
where Lemma~\ref{lem:TAIL} applies and the finitely many boundary rows
$\nu+1 \le L < \lambda_1 \le 2\nu$.
\end{enumerate}

\subsection*{C.2 The \texorpdfstring{$E = 0$}{E = 0} column, printed in full}

By Lemma~\ref{lem:TAIL}, for $L \ge \lambda_1$ a row clears in exactly $L$ steps if and only
if $s_0(\lambda) = \lambda_1$, where
$s_0(\lambda) = \mathrm{steps}([\lambda_1]^{\lambda_1+1} \cup \lambda)$ --- a quantity of
$\lambda$ alone, computable by hand with one Havel--Hakimi run. Below, every partition of
$2\nu$ is listed as $\lambda : s_0(\lambda)$, survivors ($s_0 = \lambda_1$) in bold. The
counts are $p(2\nu)$, so the lists are complete by inspection.

{\small\sloppy\noindent
$\boldsymbol{\nu = 1}$ --- $p(2) = 2$:\quad
\textbf{2:2} $\cdot$ 1+1:2

\noindent
$\boldsymbol{\nu = 2}$ --- $p(4) = 5$:\quad
\textbf{4:4} $\cdot$ 3+1:4 $\cdot$ 2+2:3 $\cdot$ 2+1+1:3 $\cdot$ 1+1+1+1:3

\noindent
$\boldsymbol{\nu = 3}$ --- $p(6) = 11$:\quad
\textbf{6:6} $\cdot$ 5+1:6 $\cdot$ 4+2:5 $\cdot$ 4+1+1:5 $\cdot$ 3+3:4 $\cdot$ 3+2+1:4
$\cdot$ 3+1+1+1:5 $\cdot$ 2+2+2:4 $\cdot$ 2+2+1+1:4 $\cdot$ 2+1+1+1+1:4 $\cdot$
1+1+1+1+1+1:4

\noindent
$\boldsymbol{\nu = 4}$ --- $p(8) = 22$:\quad
\textbf{8:8} $\cdot$ 7+1:8 $\cdot$ 6+2:7 $\cdot$ 6+1+1:7 $\cdot$ 5+3:6 $\cdot$ 5+2+1:6
$\cdot$ 5+1+1+1:7 $\cdot$ 4+4:5 $\cdot$ 4+3+1:5 $\cdot$ 4+2+2:6 $\cdot$ 4+2+1+1:6 $\cdot$
4+1+1+1+1:6 $\cdot$ 3+3+2:5 $\cdot$ 3+3+1+1:5 $\cdot$ 3+2+2+1:5 $\cdot$ 3+2+1+1+1:5 $\cdot$
3+1+1+1+1+1:6 $\cdot$ 2+2+2+2:4 $\cdot$ 2+2+2+1+1:5 $\cdot$ 2+2+1+1+1+1:5 $\cdot$
2+1+1+1+1+1+1:5 $\cdot$ 1+1+1+1+1+1+1+1:5

\noindent
$\boldsymbol{\nu = 5}$ --- $p(10) = 42$:\quad
\textbf{10:10} $\cdot$ 9+1:10 $\cdot$ 8+2:9 $\cdot$ 8+1+1:9 $\cdot$ 7+3:8 $\cdot$ 7+2+1:8
$\cdot$ 7+1+1+1:9 $\cdot$ 6+4:7 $\cdot$ 6+3+1:7 $\cdot$ 6+2+2:8 $\cdot$ 6+2+1+1:8 $\cdot$
6+1+1+1+1:8 $\cdot$ 5+5:6 $\cdot$ 5+4+1:6 $\cdot$ 5+3+2:7 $\cdot$ 5+3+1+1:7 $\cdot$
5+2+2+1:7 $\cdot$ 5+2+1+1+1:7 $\cdot$ 5+1+1+1+1+1:8 $\cdot$ 4+4+2:6 $\cdot$ 4+4+1+1:6
$\cdot$ 4+3+3:6 $\cdot$ 4+3+2+1:6 $\cdot$ 4+3+1+1+1:6 $\cdot$ 4+2+2+2:6 $\cdot$
4+2+2+1+1:7 $\cdot$ 4+2+1+1+1+1:7 $\cdot$ 4+1+1+1+1+1+1:7 $\cdot$ 3+3+3+1:5 $\cdot$
3+3+2+2:5 $\cdot$ 3+3+2+1+1:6 $\cdot$ 3+3+1+1+1+1:6 $\cdot$ 3+2+2+2+1:6 $\cdot$
3+2+2+1+1+1:6 $\cdot$ 3+2+1+1+1+1+1:6 $\cdot$ 3+1+1+1+1+1+1+1:7 $\cdot$ 2+2+2+2+2:5 $\cdot$
2+2+2+2+1+1:5 $\cdot$ 2+2+2+1+1+1+1:6 $\cdot$ 2+2+1+1+1+1+1+1:6 $\cdot$
2+1+1+1+1+1+1+1+1:6 $\cdot$ 1+1+1+1+1+1+1+1+1+1:6

\noindent
$\boldsymbol{\nu = 6}$ --- $p(12) = 77$:\quad
\textbf{12:12} $\cdot$ 11+1:12 $\cdot$ 10+2:11 $\cdot$ 10+1+1:11 $\cdot$ 9+3:10 $\cdot$
9+2+1:10 $\cdot$ 9+1+1+1:11 $\cdot$ 8+4:9 $\cdot$ 8+3+1:9 $\cdot$ 8+2+2:10 $\cdot$
8+2+1+1:10 $\cdot$ 8+1+1+1+1:10 $\cdot$ 7+5:8 $\cdot$ 7+4+1:8 $\cdot$ 7+3+2:9 $\cdot$
7+3+1+1:9 $\cdot$ 7+2+2+1:9 $\cdot$ 7+2+1+1+1:9 $\cdot$ 7+1+1+1+1+1:10 $\cdot$ 6+6:7
$\cdot$ 6+5+1:7 $\cdot$ 6+4+2:8 $\cdot$ 6+4+1+1:8 $\cdot$ 6+3+3:8 $\cdot$ 6+3+2+1:8 $\cdot$
6+3+1+1+1:8 $\cdot$ 6+2+2+2:8 $\cdot$ 6+2+2+1+1:9 $\cdot$ 6+2+1+1+1+1:9 $\cdot$
6+1+1+1+1+1+1:9 $\cdot$ 5+5+2:7 $\cdot$ 5+5+1+1:7 $\cdot$ 5+4+3:7 $\cdot$ 5+4+2+1:7 $\cdot$
5+4+1+1+1:7 $\cdot$ 5+3+3+1:7 $\cdot$ 5+3+2+2:7 $\cdot$ 5+3+2+1+1:8 $\cdot$
5+3+1+1+1+1:8 $\cdot$ 5+2+2+2+1:8 $\cdot$ 5+2+2+1+1+1:8 $\cdot$ 5+2+1+1+1+1+1:8 $\cdot$
5+1+1+1+1+1+1+1:9 $\cdot$ 4+4+4:6 $\cdot$ 4+4+3+1:6 $\cdot$ 4+4+2+2:6 $\cdot$
4+4+2+1+1:7 $\cdot$ 4+4+1+1+1+1:7 $\cdot$ 4+3+3+2:6 $\cdot$ 4+3+3+1+1:7 $\cdot$
4+3+2+2+1:7 $\cdot$ 4+3+2+1+1+1:7 $\cdot$ 4+3+1+1+1+1+1:7 $\cdot$ 4+2+2+2+2:7 $\cdot$
4+2+2+2+1+1:7 $\cdot$ 4+2+2+1+1+1+1:8 $\cdot$ 4+2+1+1+1+1+1+1:8 $\cdot$
4+1+1+1+1+1+1+1+1:8 $\cdot$ 3+3+3+3:6 $\cdot$ 3+3+3+2+1:6 $\cdot$ 3+3+3+1+1+1:6 $\cdot$
3+3+2+2+2:6 $\cdot$ 3+3+2+2+1+1:6 $\cdot$ 3+3+2+1+1+1+1:7 $\cdot$ 3+3+1+1+1+1+1+1:7
$\cdot$ 3+2+2+2+2+1:6 $\cdot$ 3+2+2+2+1+1+1:7 $\cdot$ 3+2+2+1+1+1+1+1:7 $\cdot$
3+2+1+1+1+1+1+1+1:7 $\cdot$ 3+1+1+1+1+1+1+1+1+1:8 $\cdot$ 2+2+2+2+2+2:6 $\cdot$
2+2+2+2+2+1+1:6 $\cdot$ 2+2+2+2+1+1+1+1:6 $\cdot$ 2+2+2+1+1+1+1+1+1:7 $\cdot$
2+2+1+1+1+1+1+1+1+1:7 $\cdot$ 2+1+1+1+1+1+1+1+1+1+1:7 $\cdot$
1+1+1+1+1+1+1+1+1+1+1+1:7
\par}

\medskip
\noindent
\textbf{Reading.} For every $\nu \le 6$ the only partition of $2\nu$ with
$s_0(\lambda) = \lambda_1$ is the single part $[2\nu]$. And $[2\nu]$ has unique maximum
$2\nu \ge 2$ with second-largest entry $0 \le 2\nu-2$, so Lemma~\ref{lem:FAN6p} kills it.
Hence the rows with $E = 0$ and $L \ge \lambda_1$ contribute nothing, for every $\nu \le 6$.

\subsection*{C.3 The \texorpdfstring{$E = 0$}{E = 0} boundary rows, printed in full}

Outside Lemma~\ref{lem:TAIL}'s range the rows $\nu+1 \le L < \lambda_1$ are checked
directly. They are few, and every survivor is listed with its Lemma~\ref{lem:FAN6p}
certificate, written as the pair (unique maximum $w$, second-largest entry).

\begin{center}
\begin{tabular}{@{}cccl@{}}
\toprule
$\nu$ & pairs $(L,\lambda)$ checked & survivors & certificates $(w,\text{2nd})$ \\
\midrule
$1$ & $0$ & none & --- \\
$2$ & $1$ & $(3,[4])$ & $(4,0)$ \\
$3$ & $3$ & $(5,[6])$ & $(6,0)$ \\
$4$ & $7$ & $(5,[7,1])$, $(7,[8])$ & $(7,1)$, $(8,0)$ \\
$5$ & $14$ & $(7,[9,1])$, $(9,[10])$ & $(9,1)$, $(10,0)$ \\
$6$ & $26$ & $(7,[10,1,1])$, $(9,[11,1])$, $(11,[12])$ & $(10,1)$, $(11,1)$, $(12,0)$ \\
\bottomrule
\end{tabular}
\end{center}

Every certificate has a unique maximum $w$ with second-largest entry $\le w-2$, so
Lemma~\ref{lem:FAN6p} kills every boundary survivor.

\subsection*{C.4 The \texorpdfstring{$E \ge 1$}{E >= 1} rows: a closed-form count, and every survivor}

First the count, so that the ``shapes tested'' column of Theorem~\ref{thm:GFANnu} stops
being an opaque number. For $E \ge 1$, $L$ ranges over $\nu+1 \le L \le 2\nu-E$ --- that is
$\nu-E$ values, so $E \le \nu-1$; the escape multiset $e$ is a partition of $E$ into at most
$L+1$ parts, and $L+1 \ge \nu+2 > E$, so all $p(E)$ partitions of $E$ occur; and $\lambda$
is any partition of $2\nu-E$. Hence
\[
S(\nu) \;=\; \sum_{E=1}^{\nu-1} (\nu-E)\, p(E)\, p(2\nu-E).
\]
For $\nu = 6$ this is
$5\cdot1\cdot p(11) + 4\cdot2\cdot p(10) + 3\cdot3\cdot p(9) + 2\cdot5\cdot p(8) +
1\cdot7\cdot p(7) = 280+336+270+220+105 = 1\,211$, with per-$E$ split
$\{1:280,\,2:336,\,3:270,\,4:220,\,5:105\}$ reproduced term for term by the instrumented
run; the same formula gives $0,3,24,110,397,1\,211$ for $\nu = 1,\dots,6$. A reviewing
session derived this identity independently and evaluated it from its own partition table.

Of those $S(\nu)$ shapes, the ones clearing in exactly $L$ steps number
$0,1,4,9,20,38$ for $\nu = 1,\dots,6$. Here they are, each with its Lemma~\ref{lem:FAN6p}
certificate $(w,\text{2nd})$; $e$ lists the positive escape parts only.

{\small\sloppy\noindent
$\boldsymbol{\nu = 2}$ (1):\quad
$L3\ E1\ e{=}1\ \lambda{=}3$ (3,0)

\noindent
$\boldsymbol{\nu = 3}$ (4):\quad
$L4\ E1\ e{=}1\ \lambda{=}5$ (5,0) $\cdot$
$L5\ E1\ e{=}1\ \lambda{=}5$ (5,0) $\cdot$
$L4\ E2\ e{=}1{+}1\ \lambda{=}4$ (4,0) $\cdot$
$L4\ E2\ e{=}2\ \lambda{=}3{+}1$ (3,1)

\noindent
$\boldsymbol{\nu = 4}$ (9):\quad
$L6\ E1\ e{=}1\ \lambda{=}7$ (7,0) $\cdot$
$L7\ E1\ e{=}1\ \lambda{=}7$ (7,0) $\cdot$
$L5\ E2\ e{=}1{+}1\ \lambda{=}6$ (6,0) $\cdot$
$L5\ E2\ e{=}2\ \lambda{=}5{+}1$ (5,1) $\cdot$
$L6\ E2\ e{=}1{+}1\ \lambda{=}6$ (6,0) $\cdot$
$L6\ E2\ e{=}2\ \lambda{=}5{+}1$ (5,1) $\cdot$
$L5\ E3\ e{=}1{+}1{+}1\ \lambda{=}5$ (5,0) $\cdot$
$L5\ E3\ e{=}2{+}1\ \lambda{=}4{+}1$ (4,1) $\cdot$
$L5\ E3\ e{=}3\ \lambda{=}3{+}1{+}1$ (3,1)

\noindent
$\boldsymbol{\nu = 5}$ (20):\quad
$L6\ E1\ e{=}1\ \lambda{=}8{+}1$ (8,1) $\cdot$
$L8\ E1\ e{=}1\ \lambda{=}9$ (9,0) $\cdot$
$L9\ E1\ e{=}1\ \lambda{=}9$ (9,0) $\cdot$
$L6\ E2\ e{=}2\ \lambda{=}7{+}1$ (7,1) $\cdot$
$L7\ E2\ e{=}1{+}1\ \lambda{=}8$ (8,0) $\cdot$
$L7\ E2\ e{=}2\ \lambda{=}7{+}1$ (7,1) $\cdot$
$L8\ E2\ e{=}1{+}1\ \lambda{=}8$ (8,0) $\cdot$
$L8\ E2\ e{=}2\ \lambda{=}7{+}1$ (7,1) $\cdot$
$L6\ E3\ e{=}1{+}1{+}1\ \lambda{=}7$ (7,0) $\cdot$
$L6\ E3\ e{=}2{+}1\ \lambda{=}6{+}1$ (6,1) $\cdot$
$L6\ E3\ e{=}3\ \lambda{=}5{+}1{+}1$ (5,1) $\cdot$
$L7\ E3\ e{=}1{+}1{+}1\ \lambda{=}7$ (7,0) $\cdot$
$L7\ E3\ e{=}2{+}1\ \lambda{=}6{+}1$ (6,1) $\cdot$
$L7\ E3\ e{=}3\ \lambda{=}5{+}1{+}1$ (5,1) $\cdot$
$L6\ E4\ e{=}1{+}1{+}1{+}1\ \lambda{=}6$ (6,0) $\cdot$
$L6\ E4\ e{=}2{+}1{+}1\ \lambda{=}5{+}1$ (5,1) $\cdot$
$L6\ E4\ e{=}2{+}2\ \lambda{=}4{+}2$ (4,2) $\cdot$
$L6\ E4\ e{=}2{+}2\ \lambda{=}4{+}1{+}1$ (4,1) $\cdot$
$L6\ E4\ e{=}3{+}1\ \lambda{=}4{+}1{+}1$ (4,1) $\cdot$
$L6\ E4\ e{=}4\ \lambda{=}3{+}1{+}1{+}1$ (3,1)

\noindent
$\boldsymbol{\nu = 6}$ (38):\quad
$L8\ E1\ e{=}1\ \lambda{=}10{+}1$ (10,1) $\cdot$
$L10\ E1\ e{=}1\ \lambda{=}11$ (11,0) $\cdot$
$L11\ E1\ e{=}1\ \lambda{=}11$ (11,0) $\cdot$
$L7\ E2\ e{=}1{+}1\ \lambda{=}9{+}1$ (9,1) $\cdot$
$L8\ E2\ e{=}2\ \lambda{=}9{+}1$ (9,1) $\cdot$
$L9\ E2\ e{=}1{+}1\ \lambda{=}10$ (10,0) $\cdot$
$L9\ E2\ e{=}2\ \lambda{=}9{+}1$ (9,1) $\cdot$
$L10\ E2\ e{=}1{+}1\ \lambda{=}10$ (10,0) $\cdot$
$L10\ E2\ e{=}2\ \lambda{=}9{+}1$ (9,1) $\cdot$
$L7\ E3\ e{=}2{+}1\ \lambda{=}8{+}1$ (8,1) $\cdot$
$L7\ E3\ e{=}3\ \lambda{=}7{+}1{+}1$ (7,1) $\cdot$
$L8\ E3\ e{=}1{+}1{+}1\ \lambda{=}9$ (9,0) $\cdot$
$L8\ E3\ e{=}2{+}1\ \lambda{=}8{+}1$ (8,1) $\cdot$
$L8\ E3\ e{=}3\ \lambda{=}7{+}1{+}1$ (7,1) $\cdot$
$L9\ E3\ e{=}1{+}1{+}1\ \lambda{=}9$ (9,0) $\cdot$
$L9\ E3\ e{=}2{+}1\ \lambda{=}8{+}1$ (8,1) $\cdot$
$L9\ E3\ e{=}3\ \lambda{=}7{+}1{+}1$ (7,1) $\cdot$
$L7\ E4\ e{=}1{+}1{+}1{+}1\ \lambda{=}8$ (8,0) $\cdot$
$L7\ E4\ e{=}2{+}1{+}1\ \lambda{=}7{+}1$ (7,1) $\cdot$
$L7\ E4\ e{=}2{+}2\ \lambda{=}6{+}2$ (6,2) $\cdot$
$L7\ E4\ e{=}2{+}2\ \lambda{=}6{+}1{+}1$ (6,1) $\cdot$
$L7\ E4\ e{=}3{+}1\ \lambda{=}6{+}1{+}1$ (6,1) $\cdot$
$L7\ E4\ e{=}4\ \lambda{=}5{+}1{+}1{+}1$ (5,1) $\cdot$
$L8\ E4\ e{=}1{+}1{+}1{+}1\ \lambda{=}8$ (8,0) $\cdot$
$L8\ E4\ e{=}2{+}1{+}1\ \lambda{=}7{+}1$ (7,1) $\cdot$
$L8\ E4\ e{=}2{+}2\ \lambda{=}6{+}2$ (6,2) $\cdot$
$L8\ E4\ e{=}2{+}2\ \lambda{=}6{+}1{+}1$ (6,1) $\cdot$
$L8\ E4\ e{=}3{+}1\ \lambda{=}6{+}1{+}1$ (6,1) $\cdot$
$L8\ E4\ e{=}4\ \lambda{=}5{+}1{+}1{+}1$ (5,1) $\cdot$
$L7\ E5\ e{=}1{+}1{+}1{+}1{+}1\ \lambda{=}7$ (7,0) $\cdot$
$L7\ E5\ e{=}2{+}1{+}1{+}1\ \lambda{=}6{+}1$ (6,1) $\cdot$
$L7\ E5\ e{=}2{+}2{+}1\ \lambda{=}5{+}2$ (5,2) $\cdot$
$L7\ E5\ e{=}2{+}2{+}1\ \lambda{=}5{+}1{+}1$ (5,1) $\cdot$
$L7\ E5\ e{=}3{+}1{+}1\ \lambda{=}5{+}1{+}1$ (5,1) $\cdot$
$L7\ E5\ e{=}3{+}2\ \lambda{=}4{+}2{+}1$ (4,2) $\cdot$
$L7\ E5\ e{=}3{+}2\ \lambda{=}4{+}1{+}1{+}1$ (4,1) $\cdot$
$L7\ E5\ e{=}4{+}1\ \lambda{=}4{+}1{+}1{+}1$ (4,1) $\cdot$
$L7\ E5\ e{=}5\ \lambda{=}3{+}1{+}1{+}1{+}1$ (3,1)
\par}

\medskip
\noindent
\textbf{Reading.} In every one of these $72$ rows the residue $\lambda$ has a unique maximum
$w \ge 3$ whose second-largest entry is $\le w-2$: the pairs are $(w,0)$, or $(w,1)$ with
$w \ge 3$, or $(w,2)$ with $w \ge 4$. So Lemma~\ref{lem:FAN6p} kills every $E \ge 1$
survivor, for every $\nu \le 6$. That is the ``misses = none'' column of
Theorem~\ref{thm:GFANnu}, printed as data rather than asserted as an output.

\subsection*{C.5 What a re-reader has to do}

To reproduce Theorem~\ref{thm:GFANnu} from this text it suffices to (i) recompute
$s_0(\lambda)$ for the $159$ partitions of C.2 --- one Havel--Hakimi run each, all by hand
--- and confirm that the bolded survivor is the only one at each $\nu$; (ii) check the $51$
boundary pairs of C.3; (iii) evaluate $S(\nu)$ from the closed form of C.4 and confirm that
the $72$ printed survivors are the complete survivor set for those $1\,745$ shapes; and
(iv) apply the Lemma~\ref{lem:FAN6p} certificate to each of the $9 + 72 + 6$ surviving rows.
Only step (iii)'s completeness claim still asks the reader to trust a machine run, and it is
a claim about an explicitly specified, closed-form-counted finite set rather than about an
absent file. That is the exact location of the demarcation line of
Remark~\ref{rem:demarcation}.

\begin{thebibliography}{99}

\bibitem{formalconjectures}
Google DeepMind and contributors,
\emph{formal-conjectures}: a corpus of open conjectures stated in Lean~4,
\url{https://github.com/google-deepmind/formal-conjectures}.
The statement treated here is
\texttt{FormalConjectures/WrittenOnTheWallII/GraphConjecture61.lean}, with supporting
definitions in
\texttt{FormalConjecturesForMathlib/Combinatorics/SimpleGraph/Residue.lean} and
\texttt{.../Induced.lean}. The corpus attributes the conjecture to the
\emph{Written on the Wall II} collection generated by the \texttt{Graffiti.pc} program; we
consumed the formal file and did not consult the original collection.

\bibitem{havel}
V.~Havel,
\emph{A remark on the existence of finite graphs},
\v{C}asopis P\v{e}st. Mat. \textbf{80} (1955), 477--480.
Bibliographic details as supplied by our automated literature pass.

\bibitem{hakimi}
S.~L.~Hakimi,
\emph{On realizability of a set of integers as degrees of the vertices of a linear graph},
J. Soc. Indust. Appl. Math. \textbf{10} (1962), 496--506.
Bibliographic details as supplied by our automated literature pass.

\bibitem{fms}
O.~Favaron, M.~Mah\'eo and J.-F.~Sacl\'e,
\emph{On the residue of a graph},
J. Graph Theory \textbf{15} (1991), 39--64.
The source of Fact~\ref{fact:fms}. Bibliographic details as supplied by our automated
literature pass; the inequality itself is verified computationally on the whole exhaustive
range but is not reproved here (Section~\ref{sec:limits}).

\bibitem{gk}
J.~R.~Griggs and D.~J.~Kleitman,
\emph{Independence and the Havel--Hakimi residue},
Discrete Math. \textbf{127} (1994), 209--212.
A short proof of Fact~\ref{fact:fms}. Bibliographic details as supplied by our automated
literature pass.

\bibitem{mathlib}
The mathlib Community,
\emph{The Lean mathematical library},
Proc. 9th ACM SIGPLAN Int. Conf. on Certified Programs and Proofs (CPP 2020), 367--381.
\url{https://github.com/leanprover-community/mathlib4}.
The source of the \texttt{SimpleGraph.diam} and degree-sequence definitions against which
the residue transcription of Section~\ref{sec:statement} was checked.

\bibitem{batch1}
H.~Chen,
\emph{Machine-Verified Resolutions of Four OEIS Conjectures}, 2026.
The companion report of the same pipeline on four resolved conjectures, each formalized in
Lean~4; it describes the review protocol summarised in Section~\ref{sec:conventions}.

\bibitem{fernandespaper}
H.~Chen,
\emph{The parity subgroup of $\mathrm{S}_m \times \mathrm{S}_n$ is $2$-generated: a proof of
Fernandes' conjecture}, 2026.

\bibitem{etp677}
H.~Chen,
\emph{Structure Theory of Finite 677 Magmas: a working paper}, 2026.
The companion working paper on an unfinished campaign, from which the per-assertion
verification-marker discipline of Section~\ref{sec:conventions} is taken.

\bibitem{erdosproblems}
T.~F.~Bloom (site), T.~Tao (community database) and contributors,
\emph{Erd\H{o}s Problems} and the associated AI-contributions wiki.
\url{https://www.erdosproblems.com}, \url{https://github.com/teorth/erdosproblems}.
Cited for its apparatus of caveats about selection bias and about the difference between an
argument and a verified proof, which the tiering of Section~\ref{sec:conventions} is
intended to respect.

\end{thebibliography}
